\section{Primitive-defect source audit: a proved coordinate core}
\label{sec:primitive-defect-source-audit}

This chapter records the part of the modern local source
``Primitive Defect Augmentations, Split-Zero Support, and A-Series Realizations''
which survives a line-by-line audit.  The local source was read consecutively
at its hash-pinned TeX edition (lines~1--1474).  It is an input artifact, not a
human provenance record.  The present chapter therefore cites human sources
for imported facts and labels every local construction by its actual proof
scope.  In particular, the source's final ``consolidated theorem'' is not
retained as one theorem: its summands live in different typed categories, and
several of the proposed arrows were not defined.

\subsection{Scope, source identity, and audit boundary}

The audited source has 47,978 bytes and SHA-256
\begin{quote}\small\texttt{2f3289f6a04105d3f15fe5f7c36dde286515e61f465caf1}\allowbreak\\
\texttt{582f9808e19f3d2f7}.\end{quote}
An identical TeX copy occurs in the adjacent version directory.  The source
contains the following distinct layers:

\begin{enumerate}[label=(\alph*)]
  \item literal split-zero semiring and support-linearization constructions;
  \item augmentation, Fourier, root-lattice, and Eisenstein calculations;
  \item a proposed semilocal crossed-product boundary model;
  \item principal parts of a shifted completed Hurwitz function;
  \item orbifold, Weyl-operator, CUE, spectral, F-theory, and moonshine
        interfaces; and
  \item a conjectural CUE/Epstein transform bridge.
\end{enumerate}

Only (a), the finite-dimensional portions of (b), the principal-part algebra,
and the explicitly typed finite CUE/Weyl calculations are proved below.  The
other layers are retained in the machine-readable audit as candidate
interfaces with their missing definitions and dependencies.  ``Not proved in
this chapter'' is a source-relative statement; it is not a claim about the
whole literature.

\subsection{The split-zero semiring and its two zero fibres}

Let $R$ be a nonzero commutative ring and put
\[
  G_R=R\sqcup\{\tau\}.
\]
For $r,s\in R$ define
\[
 r\oplus s=r+s,\qquad r\oplus\tau=\tau\oplus r=r,
 \qquad \tau\oplus\tau=\tau,
\]
and
\[
 r\odot s=rs,\qquad r\odot\tau=\tau\odot r=\tau,
 \qquad \tau\odot\tau=\tau.
\]

\begin{proposition}[Semiring and coordinate model]
\label{prop:primitive-split-zero}
The operations above make $G_R$ a commutative semiring with additive zero
$\tau$ and multiplicative unit $1_R$.  The map
\[
 \iota:G_R\longrightarrow D_R:=\{(0,0)\}\cup(R\times\{1\}),
 \qquad \iota(\tau)=(0,0),\quad \iota(r)=(r,1),
\]
is an injective semiring homomorphism when $D_R$ has the coordinatewise
operations inherited from $R\times\{0,1\}$ with the same split-zero tables.
The projections
\[
 p_R(\tau)=0_R,\quad p_R(r)=r,
 \qquad
 \chi_R(\tau)=0,\quad \chi_R(r)=1
\]
have fibres
\[
 p_R^{-1}(0_R)=\{\tau,0_R\},\qquad
 \chi_R^{-1}(0)=\{\tau\},\qquad
 \chi_R^{-1}(1)=R.
\]
\end{proposition}

\begin{proof}
If all entries are in $R$, associativity, commutativity, and distributivity
are those of $R$.  If an entry is $\tau$, the addition table says that it is
an identity, so every associativity case reduces to the case with the $\tau$
removed.  For multiplication, $\tau$ is absorbing, so associativity and
commutativity reduce in the same way.  For distributivity, if the multiplier
is $\tau$ both sides are $\tau$; if the multiplier and both summands are in
$R$, distributivity is inherited from $R$; and, for example,
\[
 r\odot(s\oplus\tau)=r\odot s=rs\oplus\tau
                 =(r\odot s)\oplus(r\odot\tau).
\]
The remaining cases are symmetric.  Thus the semiring axioms hold.
The displayed coordinate map is injective because $(0,0)$ is not in
$R\times\{1\}$.  Checking a pair of supported entries uses the ring laws;
checking a pair containing $\tau$ uses the two defining tables, so it
preserves both operations and the unit.  Reading the first and second
coordinates gives the three fibre identities.
\end{proof}

\begin{theorem}[Support and arithmetic prime ideals]
\label{thm:primitive-prime-ideals}
The subset $\mathfrak s_R=\{\tau\}$ is a prime ideal of $G_R$.  If $R$ is an
integral domain, then $\mathfrak z_R=\{\tau,0_R\}$ is also a prime ideal, and
\[
 \mathfrak s_R\subset\mathfrak z_R.
\]
\end{theorem}

\begin{proof}
The set $\mathfrak s_R$ is closed under addition and absorbs multiplication.
If $x\odot y=\tau$ and neither factor is $\tau$, then $x,y\in R$ and
$x\odot y=xy\in R$, a contradiction.  Hence one factor belongs to
$\mathfrak s_R$, which is primality.

Assume now that $R$ is a domain.  The set $\mathfrak z_R$ is closed under
addition and absorbs multiplication.  If $x\odot y\in\mathfrak z_R$ and
neither factor is $\tau$, then $x,y\in R$ and $xy=0_R$; the domain property
gives $x=0_R$ or $y=0_R$.  Thus one factor is in $\mathfrak z_R$.
Both ideals are proper because $R\ne\{0_R\}$, and the inclusion is immediate.
\end{proof}

The two ideals retain different information: the first separates support from
absence, while the second also identifies the supported arithmetic zero with
the semiring zero.  No identification of these ideals with objects in a
different category is made here.

\subsection{Support linearization and the $A$-series kernel}

For a finite set $I$, let $\mathcal P(I)$ be its Boolean support cube and put
\[
 \mathcal L_I(J)=\sum_{j\in J}e_j\in\mathbb Z^I,
 \qquad
 \delta_I\!\left(\sum_{i\in I}a_i e_i\right)=-\sum_{i\in I}a_i.
\]

\begin{theorem}[Exact kernel coordinates]
\label{thm:primitive-A-kernel}
If $|I|=r$, then
\[
 \ker\delta_I=\left\{(a_1,\ldots,a_r)\in\mathbb Z^r:
                    \sum_i a_i=0\right\}
\]
has basis $\alpha_i=e_i-e_{i+1}$ $(1\leq i<r)$.  Its Gram matrix for the
standard inner product is the Cartan matrix of type $A_{r-1}$; for $r=1$ the
kernel is the zero lattice $A_0$.
\end{theorem}

\begin{proof}
Choose an ordering of $I$.  Every $\alpha_i$ has coordinate sum zero.  If
$a=(a_1,\ldots,a_r)$ has sum zero, set
\[
 b_i=a_1+\cdots+a_i\qquad(1\leq i<r).
\]
The $j$th coordinate of $\sum_{i<r}b_i(e_i-e_{i+1})$ is $b_1=a_1$ for
$j=1$, is $b_j-b_{j-1}=a_j$ for $1<j<r$, and is
$-b_{r-1}=a_r$ for $j=r$.  Thus the $\alpha_i$ span the kernel.  If a linear
combination of them is zero, its first coordinate is the coefficient of
$\alpha_1$, then its second new coordinate determines the coefficient of
$\alpha_2$, and so on; hence they are independent.
Finally
\[
 \langle\alpha_i,\alpha_j\rangle=
 \begin{cases}2&i=j,\\-1&|i-j|=1,\\0&|i-j|>1,\end{cases}
\]
which is the $A_{r-1}$ Cartan matrix.  When $r=1$ there are no roots and the
kernel is $0$.
\end{proof}

The defect of a support $J$ is exactly
$\delta_I(\mathcal L_I(J))=-|J|$.  This is a statement about the displayed
free abelian group.  A Grothendieck group, representation ring, or fusion
ring may have relations and primitive objects of dimension other than one;
the slogan ``universal negative dimension'' is therefore valid only after
the explicit free, dimension-one hypotheses have been imposed.

\subsection{Cyclic Fourier projection}

Let $C_r=\langle g\mid g^r=1\rangle$ and
$\omega=e^{2\pi i/r}$.  In $\mathbb C[C_r]$ define
\[
 S_{ab}=r^{-1/2}\omega^{ab},\qquad a,b\in\mathbb Z/r\mathbb Z.
\]

\begin{proposition}[Augmentation is the trivial Fourier coordinate]
\label{prop:primitive-fourier}
For $x=\sum_{a=0}^{r-1}c_ag^a$,
\[
 \epsilon(x)=\sum_a c_a=\sqrt r\sum_a c_aS_{a0}.
\]
The nontrivial Fourier modes span
$\ker(\epsilon)\otimes_{\mathbb Z}\mathbb C$.
\end{proposition}

\begin{proof}
Since $S_{a0}=r^{-1/2}$, the first equality is immediate.  For
$b\in\mathbb Z/r\mathbb Z$ put
$f_b=\sum_a\omega^{-ab}g^a$.  The finite geometric sum gives
\[
 \sum_{a=0}^{r-1}\omega^{a(b-c)}=r\,\mathbf 1_{b=c},
\]
so the $f_b$ form a basis.  The coefficient sum of $f_b$ is $r$ for $b=0$
and $0$ otherwise.  Therefore the kernel is exactly the span of
$f_1,\ldots,f_{r-1}$.
\end{proof}

\subsection{The $C_3$ augmentation ideal and the Eisenstein Epstein sum}

For $C_3=\langle g\rangle$, the augmentation ideal has basis
$\alpha_1=1-g$, $\alpha_2=g-g^2$, with Gram matrix
\[
 \begin{pmatrix}2&-1\\-1&2\end{pmatrix}.
\]
The evaluation map $\pi_\omega:\mathbb Z[C_3]\to\mathbb Z[\omega]$,
$g\mapsto\omega$, has kernel generated by $1+g+g^2$, and
\[
 \pi_\omega(I_{C_3})=(1-\omega)\mathbb Z[\omega].
\]
Indeed, the two basis images are $1-\omega$ and
$\omega-\omega^2=\omega(1-\omega)$, and $\omega$ is a unit.

Writing $z=m+n\omega$ gives
\[
 N(z)=z\overline z=m^2-mn+n^2,
 \qquad Q_{A_2}(m,n)=2N(z).
\]

\begin{theorem}[$A_2$ Epstein zeta in its initial half-plane]
\label{thm:primitive-A2-epstein}
For $\Re(s)>1$,
\[
 Z_{A_2}(s):=\sum_{(m,n)\ne(0,0)}Q_{A_2}(m,n)^{-s}
 =2^{-s}6\,\zeta(s)L(s,\chi_{-3}).
\]
The identity then extends wherever the two cited meromorphic continuations
are defined.
\end{theorem}

\begin{proof}
The norm is multiplicative because it is the complex norm in
$K=\mathbb Q(\sqrt{-3})$.  The Eisenstein integers are Euclidean for this
norm: the triangular lattice $\mathbb Z[\omega]\subset\mathbb C$ has Voronoi
cells of circumradius $1/\sqrt3<1$, so for $z\in K$ choose
$q\in\mathbb Z[\omega]$ with $|z-q|<1$ and, after clearing a denominator,
obtain Euclidean division $a=bq+r$ with $N(r)<N(b)$.  Hence every nonzero
ideal is principal.  The six elements
\(\{\pm1,\pm\omega,\pm\omega^2\}\) are exactly the units (their norm is
one), so each nonzero principal ideal has six generators.

For $\Re(s)>1$, unique ideal factorization and multiplicativity of the ideal
norm give the Dedekind Euler product (the standard theorem is recorded in
Milne and Neukirch~\cite{milneANT2020,neukirch1999}).  Consequently
\[
 \sum_{0\ne z\in\mathbb Z[\omega]}N(z)^{-s}=6\,\zeta_K(s).
\]
The quadratic character of $K$ is $\chi_{-3}$; comparing the Euler factor at
each rational prime (split for $p\equiv1\pmod3$, inert for
$p\equiv2\pmod3$, and ramified at $3$) gives
\[
 \zeta_K(s)=\zeta(s)L(s,\chi_{-3}).
\]
Multiplying by $2^{-s}$ proves the displayed formula.  The series argument
was made only in $\Re(s)>1$; continuation is imported from the cited
Dedekind-zeta theorem rather than silently inferred from the series.
\end{proof}

\subsection{Completed Hurwitz principal parts}

The local source writes $F_t(s)=\zeta(s,1+t)$ and
\[
 \Lambda_t(s)=\pi^{-s/2}\Gamma(s/2)F_t(s).
\]
The value formula used here is DLMF~25.11.13.  The simple pole and
specialization are recorded in DLMF~25.11.1--2~\cite{apostolDLMF25}.
The following statement uses $a=1+t$ in the DLMF range $\Re(a)>0$ (or, more
generally, away from nonpositive-integer parameter singularities).

\begin{lemma}[Invariant and anti-invariant principal-part basis]
\label{lem:primitive-principal-parts}
Put $u=s(s-1)$ and $v=2s-1$.  Then
\[
 u(1-s)=u(s),\qquad v(1-s)=-v(s),\qquad v^2=1+4u,
\]
and
\[
 \frac1u=\frac1{s-1}-\frac1s,
 \qquad
 \frac vu=\frac1{s-1}+\frac1s.
\]
If $P(s)=a/s+b/(s-1)$, then
\[
 P(s)=A\frac1u+B\frac vu,\qquad
 A=\frac{b-a}{2},\quad B=\frac{a+b}{2}.
\]
\end{lemma}

\begin{proof}
The first three identities follow by direct substitution.  Since
$u=s(s-1)$ and $v=2s-1$, partial fractions give the next two identities.
Expanding $A/u+Bv/u$ gives coefficient $-A+B$ at $1/s$ and $A+B$ at
$1/(s-1)$.  Solving $a=-A+B$, $b=A+B$ gives the displayed $A,B$.
\end{proof}

\begin{theorem}[Shifted completed-Hurwitz endpoint calculation]
\label{thm:primitive-hurwitz-endpoints}
Under the parameter restriction above,
\[
 \operatorname*{Res}_{s=0}\Lambda_t=-1-2t,
 \qquad
 \operatorname*{Res}_{s=1}\Lambda_t=1,
\]
and hence
\[
 \operatorname{PP}_{\{0,1\}}(\Lambda_t)
 =(1+t)\frac1u-t\frac vu.
\]
For $t=1$ the anti-invariant term is $-v/u$, with residue $-1$ at both
endpoints.
\end{theorem}

\begin{proof}
At $s=1$, DLMF gives residue one for $\zeta(s,a)$, while
$\pi^{-1/2}\Gamma(1/2)=1$, proving the second residue.  At $s=0$,
\[
 \Gamma(s/2)=\frac2s+O(1),\qquad \pi^{-s/2}=1+O(s),
 \]
and DLMF~25.11.13 gives
$\zeta(0,1+t)=\tfrac12-(1+t)=-\tfrac12-t$.  The first residue is therefore
$2(-\tfrac12-t)=-1-2t$.  Apply Lemma~\ref{lem:primitive-principal-parts}
with $a=-1-2t$ and $b=1$; it gives $A=1+t$ and $B=-t$.  At $t=1$,
$-v/u=-1/s-1/(s-1)$, so both endpoint residues are $-1$.
\end{proof}

\begin{remark}[What the critical-line calculation actually says]
On $s=\frac12+i\gamma$,
\[
 u=-\left(\gamma^2+\frac14\right)\in\mathbb R,
 \qquad v=2i\gamma.
\]
Thus $1/u$ is real and $v/u$ is purely imaginary.  These are orthogonal only
after choosing the standard real inner product on
$\mathbb C\cong\mathbb R^2$; there is no canonical assertion of orthogonality
of complex subspaces in the source's wording.  This observation has no
implication about the location of zeros.
\end{remark}

\subsection{The finite Weyl algebra and the corrected $r=3$ matrices}

Let $V=\mathbb C^r$ with basis $|m\rangle$, $m\in\mathbb Z/r\mathbb Z$, and
define
\[
 X|m\rangle=|m+1\rangle,
 \qquad Z|m\rangle=\omega^m|m\rangle,
 \qquad W_{a,b}=X^aZ^b.
\]
Then $ZX=\omega XZ$ and
\[
 W_{a,b}W_{c,d}=\omega^{bc}W_{a+c,b+d},
 \qquad
 [W_{a,b},W_{c,d}]
   =(\omega^{bc}-\omega^{ad})W_{a+c,b+d}.
\]

\begin{theorem}[Finite Weyl basis]
\label{thm:primitive-weyl-basis}
The $r^2$ operators $W_{a,b}$ form a basis of $M_r(\mathbb C)$, and the
$r^2-1$ nonidentity operators form a basis of $\mathfrak{sl}_r(\mathbb C)$.
\end{theorem}

\begin{proof}
The relation $ZX=\omega XZ$ gives the multiplication formula, and reversing
the two factors gives the commutator.  If $a\ne0$, $X^aZ^b$ has zero diagonal;
if $a=0$ and $b\ne0$, its trace is
$\sum_{m=0}^{r-1}\omega^{bm}=0$.  Hence every nonidentity operator is
traceless.  To prove independence, suppose
$\sum_{a,b}c_{a,b}X^aZ^b=0$.  The $(m+a,m)$ entry isolates
$c_{a,b}\omega^{bm}$ as a finite Fourier polynomial in $m$; Fourier
orthogonality forces every $c_{a,b}=0$.  There are $r^2$ independent
operators in the $r^2$-dimensional matrix space, and the trace map has the
one-dimensional span of the identity as a complementary direction.  The
nonidentity span is therefore exactly the traceless subspace.
\end{proof}

For $r=3$, preserving the declared action on kets requires
\[
 X=\begin{pmatrix}0&0&1\\1&0&0\\0&1&0\end{pmatrix},
 \qquad
 X^2=\begin{pmatrix}0&1&0\\0&0&1\\1&0&0\end{pmatrix},
 \qquad
 Z=\operatorname{diag}(1,\omega,\omega^2).
\]
Direct multiplication gives $ZX=\omega XZ$.  The matrix printed as $X$ in
the audited local source is the second displayed matrix, hence is actually
$X^2$ under the source's ket convention; with that printed matrix the relation
is $ZX=\omega^2XZ$.  The correction is a relabelling of the two shifts and of
the four mixed products, not a change of the abstract Weyl algebra.  The
finite matrix calculation has no map here to a split-zero element $\tau$:
the zero matrix and unsupported absence remain distinct typed objects.

\subsection{The CUE exponent: a numerical equality and an independent block map}

Grover, Mezzadri, and Simm define, for Haar $U\in U(N)$,
\[
 \Lambda_N(z)=\det(1-zU^\dagger),\qquad
 M_{\mu,\nu}(z,N)=\mathbb E\!\left[
 \prod_i\Lambda_N^{(\mu_i)}(z)
 \prod_j\overline{\Lambda_N^{(\nu_j)}(z)}\right]
\]
and prove their microscopic equal-length asymptotic with exponent
$|\mu|+|\nu|+s^2$ under the hypotheses stated in their equal-length
microscopic theorem~\cite{groverMezzadriSimm2026}.  For $\mu=\nu=(1^s)$ this exponent is
\[
 s+s+s^2=s^2+2s=(s+1)^2-1.
\]
The last equality is a dimension count, not a statement made by the CUE
paper that its moment has a Lie-algebra-valued functor.

\begin{proposition}[Explicit block-vector-space isomorphism]
\label{prop:primitive-cue-block}
For $W=\mathbb C^s$, the map
\[
 \Theta:\operatorname{End}(W)\oplus W\oplus W^\vee
       \longrightarrow \mathfrak{sl}(W\oplus\mathbb C),\qquad
 (A,u,\varphi)\longmapsto
 \begin{pmatrix}A&u\\\varphi&-\operatorname{tr}A\end{pmatrix}
\]
is a vector-space isomorphism.
\end{proposition}

\begin{proof}
An element of $\mathfrak{sl}(W\oplus\mathbb C)$ has a unique block form
$\left(\begin{smallmatrix}B&u\\\varphi&\lambda\end{smallmatrix}\right)$ with
$\operatorname{tr}B+\lambda=0$.  It is therefore uniquely
$\Theta(B,u,\varphi)$.  Conversely every displayed block has trace zero.
Linearity and the two inverse descriptions prove bijectivity.  Transporting
the matrix commutator through this bijection gives a Lie bracket on the source
vector space, but supplies no map from the CUE moment space to that bracket.
\end{proof}

\subsection{The finite-prime boundary: local model, semilocal fibre, and exact
quotient}

This subsection closes the finite-prime local calculation and then compares it
with the two-place semilocal geometry without discarding the $p$-adic unit
coordinate.  Fix a rational prime $p$ and put $\ell=\log p$.

\begin{definition}[The one-ended valuation space]
Let
\[
 \overline{\mathbb Z}_{+}=\mathbb Z\sqcup\{\infty\}
\]
have the topology in which every integer is isolated and a neighbourhood basis
at $\infty$ is
\[
 N_M=\{\infty\}\cup\{m\in\mathbb Z:m\geq M\},\qquad M\in\mathbb Z.
\]
Thus only the positive end is compactified.  Define
\[
 \bar v_p:\mathbb Q_p\longrightarrow\overline{\mathbb Z}_{+},\qquad
 \bar v_p(x)=
 \begin{cases}v_p(x),&x\ne0,\\ \infty,&x=0.
 \end{cases}
\]
\end{definition}

\begin{lemma}[Exact valuation quotient]
\label{lem:primitive-valuation-quotient}
The map $\bar v_p$ is the topological quotient by multiplication by
$\mathbb Z_p^\times$.  In particular,
\[
 \mathbb Q_p/\mathbb Z_p^\times\ \xrightarrow{\ \sim\ }\
 \overline{\mathbb Z}_{+}
\]
is a homeomorphism, not merely a bijection of strata.
\end{lemma}

\begin{proof}
For $m\in\mathbb Z$ the fibre is
$\bar v_p^{-1}(m)=p^m\mathbb Z_p^\times$, one
$\mathbb Z_p^\times$-orbit, while the fibre over $\infty$ is $\{0\}$.
The inverse images
\[
 \bar v_p^{-1}(\{m\})=p^m\mathbb Z_p^\times,
 \qquad
 \bar v_p^{-1}(N_M)=p^M\mathbb Z_p
\]
are open, so $\bar v_p$ is continuous.  It is also open.  An open set
containing a nonzero $x$ contains a sufficiently small ball on which the
valuation is constant, hence its image contains the isolated point $v_p(x)$.
An open set containing $0$ contains $p^M\mathbb Z_p$ for some $M$, hence its
image contains $N_M$.  A continuous open surjection is a quotient map, and its
fibres are exactly the stated unit orbits.
\end{proof}

Set
\[
 Y_\ell=\overline{\mathbb Z}_{+}\times\mathbb R,
 \qquad
 T^n(m,y)=(m+n,y+n\ell),
 \qquad
 T^n(\infty,y)=(\infty,y+n\ell).
\]
Translation by an integer preserves the positive-end neighbourhoods $N_M$,
so every $T^n$ is a homeomorphism.  The space $Y_\ell$ is locally compact and
Hausdorff: an integer has a compact singleton neighbourhood, and
$N_M\times[a,b]$ is a compact neighbourhood of each $(\infty,y)$ after
shrinking the interval.

\begin{proposition}[Orbit coordinates and the spiral compactification]
\label{prop:primitive-local-orbit-space}
The $\mathbb Z$-action on $Y_\ell$ is free and proper.  Its orbit space is
\[
 X_\ell=\mathbb R\sqcup S^1_\ell,
 \qquad S^1_\ell=\mathbb R/\ell\mathbb Z,
\]
under the orbit map
\[
 q_\ell(m,y)=y-m\ell,\qquad
 q_\ell(\infty,y)=[y]_\ell.
\]
The topology is determined exactly as follows: the two displayed pieces have
their usual topologies, and a generic net $t_j\in\mathbb R$ converges to
$[a]_\ell\in S^1_\ell$ precisely when
\[
 t_j\longrightarrow-\infty,\qquad [t_j]_\ell\longrightarrow[a]_\ell.
\]
Equivalently, $X_\ell$ is homeomorphic to the closed subset of $\mathbb C$
formed by the unit circle and the disjoint spiral
\[
 t\longmapsto (1+e^t)\exp(2\pi i t/\ell),
\]
with $[a]_\ell$ sent to $\exp(2\pi i a/\ell)$.
\end{proposition}

\begin{proof}
If $T^n(m,y)=(m,y)$, then $m+n=m$; at the boundary,
$y+n\ell=y$.  In either case $n=0$, so the action is free.  If compact sets
$K,L\subset Y_\ell$ satisfy $T^nK\cap L\ne\varnothing$, the projections of
$K$ and $L$ to the real coordinate are bounded, while that coordinate changes
by $n\ell$.  Only finitely many $n$ can occur, which is properness for this
discrete action.

On the generic part, $t=y-m\ell$ is invariant and every orbit meets
$\{0\}\times\mathbb R$ once.  On the boundary the quotient is
$\mathbb R/\ell\mathbb Z$.  For the cross-stratum topology, representatives
$(m_j,y_j)$ converge to $(\infty,a)$ exactly when $m_j\to+\infty$ and
$y_j\to a$.  Their invariant coordinates then satisfy
$t_j=y_j-m_j\ell\to-\infty$ and $[t_j]_\ell=[y_j]_\ell\to[a]_\ell$.
Conversely, from those two conditions choose integers $m_j$ such that
$y_j=t_j+m_j\ell\to a$; necessarily $m_j\to+\infty$, giving representatives
converging to $(\infty,a)$.  The spiral has radius $1+e^t$, so it is injective;
the same convergence criterion identifies its complete finite accumulation
set with the unit circle.  This proves the last description and the asserted
homeomorphism.
\end{proof}

Let $U=\mathbb Z\times\mathbb R\subset Y_\ell$ and
$F=\{\infty\}\times\mathbb R$.  Since $\mathbb Z$ is amenable, full and
reduced crossed products agree and the invariant open--closed decomposition
gives
\[
 0\longrightarrow I_p\longrightarrow E_p\longrightarrow Q_p\longrightarrow0,
\]
where
\[
 I_p=C_0(U)\rtimes\mathbb Z,\qquad
 E_p=C_0(Y_\ell)\rtimes\mathbb Z,\qquad
 Q_p=C_0(F)\rtimes\mathbb Z.
\]
For a free proper action, the Green--Rieffel module identifies the crossed
product with the compact operators on the quotient module
\cite[Theorem~1.3.22]{meslandSengun2025}.  Applied to $Y_\ell$, $U$, and $F$,
it gives an extension-level Morita equivalence
\[
\begin{array}{ccccccccc}
0&\to&I_p&\to&E_p&\to&Q_p&\to&0\\
 &&\sim_M\downarrow&&\sim_M\downarrow&&\sim_M\downarrow\\
0&\to&C_0(\mathbb R)&\to&C_0(X_\ell)&\to&C(S^1_\ell)&\to&0.
\end{array}
\]
Here extension-level compatibility is not an extra assumption.  In the
standard module obtained by completing $C_c(Y_\ell)$, the closed submodule
generated by $C_0(\mathbb R)\subset C_0(X_\ell)$ is the completion of
$C_c(U)$, and its quotient is the completion of $C_c(F)$.  The corresponding
linking-algebra ideal and quotient therefore give the displayed commuting
diagram.  Functoriality of the six-term sequence transports its connecting
maps.

\begin{theorem}[The local finite-prime boundary coefficient]
\label{thm:primitive-local-boundary}
Choose the generator of $K_1(C_0(\mathbb R))$ represented by
\[
 v(t)=\frac{t-i}{t+i};
\]
as $t$ increases, this loop winds counterclockwise.  Under the explicit Green
modules above, the connecting map satisfies
\[
 \partial_p:K_0(Q_p)\longrightarrow K_1(I_p),
 \qquad \partial_p([1])=-[v].
\]
\end{theorem}

\begin{proof}
Define a function on the entire quotient $X_\ell$ by
\[
 h(t)=\phi(t)=\frac{1}{1+e^t}\quad(t\in\mathbb R),
 \qquad h([a]_\ell)=1\quad([a]_\ell\in S^1_\ell).
\]
The convergence criterion in Proposition~\ref{prop:primitive-local-orbit-space}
gives continuity at every boundary phase because $\phi(t)\to1$ as
$t\to-\infty$.  The only escaping end of $X_\ell$ is $t\to+\infty$, where
$\phi(t)\to0$, so $h\in C_0(X_\ell)$.  Its quotient is the constant projection
$1\in C(S^1_\ell)$.

Blackadar's explicit exponential formula for the six-term connecting map
\cite[Section~9.3.2]{blackadar1986} therefore gives
\[
 \partial[1]=[u],\qquad
 u(t)=\exp\!\left(2\pi i\phi(t)\right)\in 1+C_0(\mathbb R).
\]
Both endpoint limits of $u$ are $1$.  Along the oriented compactification of
$\mathbb R$, its continuous argument $2\pi\phi(t)$ decreases from $2\pi$ to
$0$, hence its winding number is $-1$.  Thus $[u]=-[v]$.  Naturality for the
extension-level Green modules gives the identical coefficient in the crossed
product extension.
\end{proof}

We now retain the unit coordinate which the local model suppresses.  The
two-place model of Connes, Consani, and Marcolli is
$(\mathbb Q_p\times\mathbb R)/(\{\pm1\}\times p^{\mathbb Z})$; its
$p$-boundary is the circle $\mathbb R_+^\times/p^{\mathbb Z}$, and its
logarithmic scale is $\mathbb R/\ell\mathbb Z$
\cite[Section~8.2]{connesConsaniMarcolli2007}.  Restrict to the generic stratum
and this $p$-boundary, namely
\[
 W_p=\mathbb Q_p\times\mathbb R^\times,
 \qquad G_p=\{\pm p^n:n\in\mathbb Z\},
\]
with diagonal multiplication.  This action is free because the real coordinate
is nonzero.  It is proper because boundedness of the absolute value of that
coordinate in two compact sets bounds $n$, leaving only the finite sign.

\begin{proposition}[Complete semilocal orbit coordinates]
\label{prop:primitive-semilocal-orbits}
The orbit space $\mathcal S_p=W_p/G_p$ is
\[
 \mathcal S_p=(\mathbb Z_p^\times\times\mathbb R)\sqcup S^1_\ell.
\]
After using the sign to choose $y>0$, its coordinates are
\[
 (x,y)\longmapsto
 \left(p^{-m}x,\,\log y-m\ell\right),
 \qquad m=v_p(x),\quad x\ne0,
\]
and
\[
 (0,y)\longmapsto[\log y]_\ell.
\]
A generic net $(a_j,t_j)$ converges to $[b]_\ell$ exactly when
\[
 t_j\to-\infty,\qquad [t_j]_\ell\to[b]_\ell;
\]
there is no condition on $a_j\in\mathbb Z_p^\times$.
\end{proposition}

\begin{proof}
The sign choice is unique because $y\ne0$.  For $x\ne0$, write uniquely
$x=p^m a$ with $m=v_p(x)$ and $a\in\mathbb Z_p^\times$.  Multiplication by
$p^n$ changes $(m,\log y)$ to $(m+n,\log y+n\ell)$ and leaves both $a$ and
$t=\log y-m\ell$ invariant.  At $x=0$, the remaining orbit coordinate is
$[\log y]_\ell$.  These formulas give a bijection.

For convergence to $(0,y)$ choose orbit representatives
$(p^{m_j}a_j,y_j)$ with $y_j\to y$.  They approach the boundary precisely when
$m_j\to+\infty$; this convergence is uniform in the unit $a_j$ because
$|p^{m_j}a_j|_p=p^{-m_j}$.  Consequently
$t_j=\log y_j-m_j\ell\to-\infty$ and
$[t_j]_\ell=[\log y_j]_\ell\to[\log y]_\ell$.  The converse follows by
choosing integers $m_j\to+\infty$ for which
$t_j+m_j\ell\to b$ and taking representatives
$(p^{m_j}a_j,\exp(t_j+m_j\ell))$.  This also proves the topological statement.
\end{proof}

The compact group $K_p=\mathbb Z_p^\times$ acts on $\mathcal S_p$ by
multiplication in the first generic coordinate and fixes the boundary circle.  Its
orbit space is exactly $X_\ell$.  Write the quotient map as
\[
 \Pi_p:\mathcal S_p\longrightarrow X_\ell,\qquad
 \Pi_p(a,t)=t,\qquad \Pi_p([b]_\ell)=[b]_\ell.
\]
Thus $\Pi_p^{-1}(t)=\mathbb Z_p^\times$ on the generic stratum and
$\Pi_p^{-1}([b]_\ell)=\{[b]_\ell\}$ on the boundary.  Pullback gives the exact
commutative diagram
\[
\begin{array}{ccccccccc}
0&\to&C_0(\mathbb R)&\to&C_0(X_\ell)&\to&C(S^1_\ell)&\to&0\\
 &&\downarrow\scriptstyle{\pi^*}&&\downarrow\scriptstyle{\Pi_p^*}
 &&\Vert\\
0&\to&C_0(\mathbb Z_p^\times\times\mathbb R)&\to&C_0(\mathcal S_p)
 &\to&C(S^1_\ell)&\to&0,
\end{array}
\]
where $\pi(a,t)=t$.  The lower extension is Morita equivalent, by the same
restricted Green-module argument, to the natural semilocal crossed-product
extension
\[
\begin{split}
0\to {}&C_0(\mathbb Q_p^\times\times\mathbb R^\times)\rtimes G_p
\longrightarrow C_0(W_p)\rtimes G_p\\
&\longrightarrow C_0(\{0\}\times\mathbb R^\times)\rtimes G_p\to0.
\end{split}
\]

\begin{corollary}[The semilocal boundary class and its exact relation to the
local coefficient]
\label{cor:primitive-semilocal-boundary}
Under the suspension/Bott identification
\[
 K_1\!\left(C_0(\mathbb Z_p^\times\times\mathbb R)\right)
 \cong K_0(C(\mathbb Z_p^\times))
 \cong C(\mathbb Z_p^\times,\mathbb Z),
\]
the semilocal connecting map sends
\[
 [1]\longmapsto \bigl(a\mapsto-1\bigr).
\]
It is the pullback of the local class $-1$ along $\pi$.
\end{corollary}

\begin{proof}
Naturality of the displayed extension diagram gives
$\partial_{\rm sl}=\pi^*_*\partial_p$.  Equivalently, the same lift is
\[
 h(a,t)=\frac1{1+e^t},\qquad h([b]_\ell)=1,
\]
and its exponential is the clockwise unitary from
Theorem~\ref{thm:primitive-local-boundary}, constant in $a$.  Since
$\mathbb Z_p^\times$ is compact and totally disconnected, stable projection
classes in $C(\mathbb Z_p^\times)$ are their locally constant integer-valued
rank functions.  Pointwise winding is $-1$ at every $a$, proving the formula.
\end{proof}

\begin{proposition}[Why the unit fibre cannot be removed by Morita
equivalence]
\label{prop:primitive-semilocal-morita-obstruction}
The map $\Pi_p$ is a proper quotient with the fibres displayed above, but it is
not a homeomorphism.  Moreover, the natural semilocal crossed-product algebra
$C_0(W_p)\rtimes G_p$ is not strongly Morita equivalent to $E_p$.
\end{proposition}

\begin{proof}
The compact group $K_p$ acts continuously on $\mathcal S_p$, so its orbit map
$\Pi_p$ is closed and proper.  On the generic open stratum it forgets the
entire $K_p$ coordinate; for example, a continuous function supported in a
proper clopen subset of $K_p$ does not lie in the image of $\Pi_p^*$.

For a topological obstruction independent of this particular formula, note
that $K_p$ is totally disconnected and has no isolated points.  Hence no point
of $K_p\times\mathbb R$ is locally connected: the projection of a connected
set to $K_p$ is a singleton, but $\{a\}\times\mathbb R$ has empty interior.
The generic stratum is dense in $\mathcal S_p$, so $\mathcal S_p$ has no
nonempty open set consisting of locally connected points.  By contrast, the
generic stratum $\mathbb R\subset X_\ell$ is a nonempty open set of locally
connected points.  Thus the two orbit spaces are not homeomorphic.

Both group actions are free and proper, so their crossed products are strongly
Morita equivalent to $C_0(\mathcal S_p)$ and $C_0(X_\ell)$, respectively
\cite[Theorem~1.3.22]{meslandSengun2025}.  Strong Morita equivalence of
commutative $C^*$-algebras induces a homeomorphism of their primitive spectra
\cite[p.~208]{green1978}.
The proved topological obstruction and transitivity of Morita equivalence rule
out a Morita equivalence between the two crossed products.
\end{proof}

Consequently the local coefficient $-1$ is correct, and its natural semilocal
image is the constant class $a\mapsto-1$.  What fails is the audited source's
claim that forgetting the $\mathbb Z_p^\times$ coordinate is itself a Morita
reduction.  A different spherical corner or unit-character projection could
still define another map, but it must be stated and proved; it is not supplied
by the source or by the two-place construction.  At this stage no larger
multi-place algebra follows from the two-place theorem alone; Subsection
\ref{subsec:unstabilized-semilocal-one-zero} constructs the natural finite-set
candidate independently and records exactly what it retains and forgets.

\paragraph{Executable boundary check.}
The bounded replay is
\begin{center}
\footnotesize
\path{certificates/globalization_nte_finite_prime_boundary/verify_finite_prime_boundary.py}.
\end{center}
It checks the two invariant coordinates, the
spiral radius limits, the derivative and endpoint values of $\phi$, the
clockwise winding $-1$, the counterclockwise Cayley generator $+1$, and the
pointwise-constant semilocal winding.  It does not certify the quotient
topologies, the Green--Rieffel theorem, extension-level Morita compatibility,
or the primitive-spectrum obstruction; those are the standard proofs printed
above.

\subsection{The archimedean reflection extension: exact matrices, boundary
map, and domain-sensitive index comparison}

The next source locus, lines~573--624, defines one genuine local object:
the reflection crossed product $C_0(\mathbb R)\rtimes(\mathbb Z/2)$.  It does
not define the larger multi-place algebra mentioned above.  Lines~626--648
then write a scalar ``global component formula'', but do not provide a global
algebra, ideal, exact sequence, codiagonal map, or component-to-global
functor.  We therefore first compute the local crossed product exactly, then
separate three inequivalent domains for the displayed differential
expression, and only after that type the additional map required by the
global formula.

Let $\sigma(x)=-x$, let $u$ denote the unitary implementing $\sigma$, and put
\[
 S=\begin{pmatrix}0&1\\1&0\end{pmatrix},\qquad
 H=\frac1{\sqrt2}\begin{pmatrix}1&1\\1&-1\end{pmatrix},\qquad
 D_2=\left\{\begin{pmatrix}z_+&0\\0&z_-\end{pmatrix}:z_\pm\in\mathbb C\right\}.
\]
The matrix $H$ is retained: it is the explicit change from the orbit basis
$(r,-r)$ to the two character coordinates of $\mathbb Z/2$.

\begin{theorem}[Reflection crossed-product boundary]
\label{thm:primitive-archimedean-boundary}
There is a $*$-isomorphism
\[
 \Phi:C_0(\mathbb R)\rtimes_\sigma(\mathbb Z/2)
 \longrightarrow
 \mathcal A_\infty
 :=\{G\in C_0([0,\infty),M_2(\mathbb C)):G(0)\in D_2\}
\]
whose value on $f_0+f_1u$ is
\begin{equation}
 \Phi(f_0+f_1u)(r)
 =H\begin{pmatrix}
       f_0(r)&f_1(r)\\
       f_1(-r)&f_0(-r)
    \end{pmatrix}H,\qquad r\geq0.
 \label{eq:primitive-archimedean-Phi}
\end{equation}
Evaluation at the fixed point gives the exact sequence
\begin{equation}
 0\longrightarrow C_0((0,\infty),M_2)
 \longrightarrow\mathcal A_\infty
 \xrightarrow{\operatorname{ev}_0}D_2\longrightarrow0.
 \label{eq:primitive-archimedean-extension}
\end{equation}
Under $D_2\cong\mathbb C\oplus\mathbb C$, the projections
\[
 p_+=\frac{1+u}{2},\qquad p_-=\frac{1-u}{2}
\]
map to $e_{11}$ and $e_{22}$, respectively.

Identify
\[
 K_1(C_0((0,\infty),M_2))\cong\mathbb Z
\]
by declaring the loop
\[
 \beta(r)=\exp\!\left(2\pi i\frac{r}{1+r}\right),\qquad
 r\in[0,\infty],
\]
to have winding $+1$.  Then the exponential connecting map of
\eqref{eq:primitive-archimedean-extension} is
\begin{equation}
 \partial_\infty:\mathbb Z[p_+]\oplus\mathbb Z[p_-]\longrightarrow\mathbb Z,
 \qquad
 \partial_\infty(a,b)=-(a+b).
 \label{eq:primitive-archimedean-boundary-map}
\end{equation}
Consequently
\[
 \partial_\infty[p_+]=\partial_\infty[p_-]=-1,\qquad
 \partial_\infty([p_+]+[p_-])=-2,\qquad
 \ker\partial_\infty=\mathbb Z(1,-1).
\]
\end{theorem}

\begin{proof}
In the orbit basis of $\{r,-r\}$, the covariant matrix of the algebraic
crossed-product element $f_0+f_1u$ is
\[
 M_f(r)=
 \begin{pmatrix}
   f_0(r)&f_1(r)\\
   f_1(-r)&f_0(-r)
 \end{pmatrix}.
\]
At $r=0$ it has the form $\bigl(\begin{smallmatrix}a&b\\b&a\end{smallmatrix}\bigr)$,
and direct multiplication gives
\[
 H\begin{pmatrix}a&b\\b&a\end{pmatrix}H
 =\begin{pmatrix}a+b&0\\0&a-b\end{pmatrix}.
\]
Thus \eqref{eq:primitive-archimedean-Phi} lands in $\mathcal A_\infty$.

Conversely, for $G\in\mathcal A_\infty$ put $M(r)=HG(r)H$.  For $r>0$
define, without identifying any of the four coordinates,
\begin{equation}
 \begin{aligned}
 f_0(r)&=M_{11}(r),& f_0(-r)&=M_{22}(r),\\
 f_1(r)&=M_{12}(r),& f_1(-r)&=M_{21}(r).
 \end{aligned}
 \label{eq:primitive-archimedean-inverse-coordinates}
\end{equation}
Because $G(0)$ is diagonal, $M(0)$ has the form
$\bigl(\begin{smallmatrix}a&b\\b&a\end{smallmatrix}\bigr)$.  Hence the two
definitions of $f_0(0)$ agree, the two definitions of $f_1(0)$ agree, and all
four half-line coordinates glue continuously.  Since $G$ vanishes at
infinity, $f_0,f_1\in C_0(\mathbb R)$.  Formula
\eqref{eq:primitive-archimedean-inverse-coordinates} is therefore an explicit
inverse to $\Phi$.

For completeness at the norm level, Mesland--\c Seng\"un define the full and
reduced crossed-product norms and the regular Hilbert-module representation
in their Sections~1.3.16--1.3.20
\cite[pp.~12--13]{meslandSengun2025}.  The group $\mathbb Z/2$ is amenable by
the explicit invariant mean $m(h)=\frac12(h(0)+h(1))$.  Applied to the
semidirect-product Fell bundle of this action, Exel's approximation theorem
therefore identifies the full and reduced cross-sectional algebras
\cite[Theorems~20.6--20.7, pp.~148--149]{exel2017}.  The regular norm here is
precisely $\sup_{r\geq0}\|M_f(r)\|$.  Thus the mutually inverse coordinate maps
above extend isometrically to the stated $C^*$-algebras.

Evaluation at zero is onto $D_2$: for any $(z_+,z_-)$, choose
$h\in C_0([0,\infty))$ with $h(0)=1$ and use
$G(r)=h(r)\operatorname{diag}(z_+,z_-)$.  Its kernel consists exactly of the
matrix functions vanishing at both $0$ and infinity, namely
$C_0((0,\infty),M_2)$.  Moreover
\[
 HSH=\operatorname{diag}(1,-1),
\]
so $p_+$ and $p_-$ become $e_{11}$ and $e_{22}$.

Take the two self-adjoint lifts
\[
 x_+(r)=\begin{pmatrix}(1+r)^{-1}&0\\0&0\end{pmatrix},\qquad
 x_-(r)=\begin{pmatrix}0&0\\0&(1+r)^{-1}\end{pmatrix}.
\]
The exponential boundary formula gives
\[
 \partial_\infty[p_\pm]=[\exp(2\pi ix_\pm)]
\]
\cite[Theorem~9.3.1 and Section~9.3.2]{blackadar1986}.  In either case the
determinant of the nontrivial block is
\[
 \exp\!\left(\frac{2\pi i}{1+r}\right)
 =\exp\!\left(-2\pi i\frac r{1+r}\right)=\beta(r)^{-1}.
\]
It starts and ends at $1$ and winds once clockwise, hence represents $-1$ in
the declared generator.  Additivity gives
\eqref{eq:primitive-archimedean-boundary-map}, including its displayed
kernel.  No sign has been normalized away.
\end{proof}

The boundary computation is therefore correct.  The source's subsequent
index comparison requires a separate domain audit.

\begin{proposition}[The displayed differential expression does not determine
an APS index]
\label{prop:primitive-archimedean-domain-audit}
Let
\[
 F_t=(1-2t)I_2,\qquad 0\leq t\leq1,\qquad
 \mathcal D_F=\frac d{dt}+F_t.
\]
Then the following three exact realizations must not be conflated.
\begin{enumerate}[label=(\alph*)]
 \item Extend $F$ continuously by $+I_2$ for $t\leq0$ and by $-I_2$ for
 $t\geq1$, and let
 \[
 D_{\mathbb R}:W^{1,2}(\mathbb R,\mathbb C^2)
 \longrightarrow L^2(\mathbb R,\mathbb C^2),
 \qquad D_{\mathbb R}=\partial_t+F_t.
 \]
 Then $\operatorname{Ind}(D_{\mathbb R})=-2$.
 \item On the compact interval, with domain
 $H^1([0,1],\mathbb C^2)$ and no endpoint condition,
 $\operatorname{Ind}(\mathcal D_F)=+2$.
 \item On the compact interval, with domain
 $\{v\in H^1:v(0)=v(1)=0\}$,
 $\operatorname{Ind}(\mathcal D_F)=-2$.
\end{enumerate}
Hence the local source's notation
$\operatorname{Ind}_{\mathrm{APS}}(\mathcal D_F)=2$ is underdetermined until
its endpoint projectors, operator domain, adjoint domain, and orientation
convention are supplied.  For realization~(a), the exact relation is
\[
 \partial_\infty([p_+]+[p_-])
 =\operatorname{Ind}(D_{\mathbb R})=-2,
\]
whereas realization~(b) gives
$\partial_\infty([p_+]+[p_-])=-\operatorname{Ind}(\mathcal D_F)$.
\end{proposition}

\begin{proof}
Doll, Schulz-Baldes, and Waterstraat count negative-to-positive crossings
positively and prove, for
\[
 D_H=\partial_t-H_t:W^{1,2}(\mathbb R,\mathbb C^N)\to
 L^2(\mathbb R,\mathbb C^N),
\]
that
\begin{equation}
 \operatorname{Ind}(D_H)=-\operatorname{Sf}(H)
 \label{eq:primitive-doll-index-spectral-flow}
\end{equation}
when the asymptotic limits are invertible
\cite[Definition~1.1.3 and Theorem~3.5.1, pp.~3 and 78--82]
{dollSchulzBaldesWaterstraat2023}.  Their complete proof identifies the kernel
and cokernel through the stable and unstable spaces and obtains
\[
 \operatorname{Ind}(D_H)
 =\dim E^u(H_{-\infty})-\dim E^u(H_{+\infty}).
\]

Here $D_{\mathbb R}=D_H$ for $H=-F$.  The path $F$ goes from $+I_2$ to
$-I_2$, so the retained signature formula gives
\[
 \operatorname{Sf}(F)=\frac12((-2)-2)=-2,
 \qquad
 \operatorname{Sf}(-F)=\frac12(2-(-2))=2.
\]
Equation~\eqref{eq:primitive-doll-index-spectral-flow} therefore yields
$\operatorname{Ind}(D_{\mathbb R})=-2$.

On $[0,1]$, every homogeneous solution is
\[
 v(t)=e^{-t+t^2}c,\qquad c\in\mathbb C^2.
\]
With no endpoint condition this gives a two-dimensional kernel.  Variation of
constants with initial value $v(0)=0$ solves $\mathcal D_Fv=g$ for every
$g\in L^2$, so the operator is onto and its index is $2$.

With both endpoint values zero, the nonvanishing scalar
$e^{-t+t^2}$ forces $c=0$, hence the kernel is zero.  Integration by parts
shows that the adjoint is $-\partial_t+F_t$ with no endpoint condition.  Its
homogeneous solutions are
\[
 w(t)=e^{t-t^2}d,\qquad d\in\mathbb C^2,
\]
so the cokernel has dimension two and the index is $-2$.  These two compact
interval domains and the whole-line Fredholm domain are distinct data.  Since
the source supplies none of the boundary projections or domains needed to
select an APS realization, its bare differential expression cannot select
one of the resulting indices.
\end{proof}

\begin{proposition}[The scalar three-component formula needs an additional
codiagonal]
\label{prop:primitive-global-codiagonal}
Let the already proved finite-prime boundary be
$\partial_p(k)=-k$, and let $\partial_\infty$ be
\eqref{eq:primitive-archimedean-boundary-map}.  The direct sum of the two
defined extensions has connecting map
\[
 \partial_p\oplus\partial_\infty:
 \mathbb Z^3\longrightarrow\mathbb Z^2,\qquad
 (k,n_+,n_-)\longmapsto(-k,-n_+-n_-).
\]
The source's scalar formula is not this direct-sum boundary.  It is the
composite with the additional codiagonal
\[
 \Sigma:\mathbb Z^2\to\mathbb Z,\qquad \Sigma(a,b)=a+b,
\]
namely
\[
 (\Sigma\circ(\partial_p\oplus\partial_\infty))(k,n_+,n_-)
 =-(k+n_++n_-).
\]
Its kernel has the explicit $A_2$ basis
\[
 b_1=(1,-1,0),\qquad b_2=(0,1,-1),\qquad
 (\langle b_i,b_j\rangle)_{i,j}
 =\begin{pmatrix}2&-1\\-1&2\end{pmatrix}.
\]
By contrast, the direct-sum boundary has the rank-one kernel
$\mathbb Z(0,1,-1)$.  Thus the $A_2$ calculation is a correct calculation for
the displayed scalar composite.  At this stage of the source audit it is not
yet the boundary map of a defined multi-place $C^*$-extension: that admission
requires an algebraic construction inducing $\Sigma$, and the audited source
supplies no such map.  The independent stable Busby construction in
Theorem~\ref{thm:global-place-stable-extension} below closes this construction
obligation after stabilization; it does not retroactively attribute the map to
the local source or identify the source's undefined unstabilized global algebra.
\end{proposition}

\begin{proof}
The direct-sum formula follows component by component from the two proved
connecting maps.  Applying $\Sigma$ gives the displayed row vector
$-(1,1,1)$.  Both $b_1$ and $b_2$ lie in its kernel, they are independent, and
their dot products give the displayed Cartan matrix.  The direct-sum map has
two independent rows, so its kernel is exactly
$\mathbb Z(0,1,-1)$.  This proves both kernel statements and isolates the
additional information carried by $\Sigma$.
\end{proof}

\paragraph{Executable archimedean check.}
The bounded replay is
\begin{center}
\footnotesize
\path{certificates/globalization_nte_archimedean_boundary/verify_archimedean_boundary.py}.
\end{center}
It checks the Hadamard coordinate change, fixed-point boundary form, character
projections, explicit lifts and winding, both spectral-flow signs, the three
domain-sensitive index counts, the direct-sum and codiagonal matrices, and the
$A_2$ Gram matrix.  It does not certify the crossed-product completion, the
six-term theorem, Fredholmness, or an APS boundary theorem; the exact standard
proofs and source dependencies are printed above.


\subsection{Global place balancing: an exact lattice and a stable extension
theorem}

The preceding local calculations determine all coefficients that occur in the
source's scalar component formula, but they do not yet make that formula a
connecting map of one global extension.  This distinction produces a concrete
construction problem rather than a slogan.

Fix an ordered list of distinct primes
\[
 P=(p_1,\ldots,p_n),\qquad n\geq0,
\]
and retain the two archimedean character coordinates $p_+$ and $p_-$.  In the
ordered coordinate basis
\[
 ( [1]_{p_1},\ldots,[1]_{p_n},[p_+],[p_-])
\]
define the already determined direct-sum homomorphism
\begin{equation}
 \Delta_P:\mathbb Z^{n+2}\longrightarrow\mathbb Z^{n+1},
 \qquad
 \Delta_P(k_1,\ldots,k_n,n_+,n_-)
 =(-k_1,\ldots,-k_n,-n_+-n_-).
 \label{eq:global-place-direct-sum-boundary}
\end{equation}
Its first $n$ coordinates are the finite-prime coefficients proved in
Theorem~\ref{thm:primitive-local-boundary}; its last coordinate is the
archimedean map proved in
Theorem~\ref{thm:primitive-archimedean-boundary}.  No place, sign, or
character coordinate has been identified in writing
\eqref{eq:global-place-direct-sum-boundary}.

Let
\begin{equation}
 \Sigma_P:\mathbb Z^{n+1}\longrightarrow\mathbb Z,
 \qquad
 \Sigma_P(y_1,\ldots,y_n,y_\infty)
 =\sum_{i=1}^{n}y_i+y_\infty
 \label{eq:global-place-codiagonal}
\end{equation}
and put
\begin{equation}
 \delta_P^{\mathrm{bal}}=\Sigma_P\circ\Delta_P,
 \qquad
 \delta_P^{\mathrm{bal}}(k_1,\ldots,k_n,n_+,n_-)
 =-\left(\sum_{i=1}^{n}k_i+n_++n_-\right).
 \label{eq:global-place-balancing-row}
\end{equation}
At this point $\delta_P^{\mathrm{bal}}$ is a completely defined homomorphism of
integer lattices.  It is not yet declared to be the boundary map of a global
$C^*$-extension.

\begin{theorem}[Exact balancing lattice]
\label{thm:global-place-balancing-lattice}
Write $m=n+2$ and let $e_1,\ldots,e_m$ be the retained standard basis of
$\mathbb Z^m$.  Then $\delta_P^{\mathrm{bal}}$ is onto and
\[
 \ker\delta_P^{\mathrm{bal}}
 =\left\{x\in\mathbb Z^m:\sum_{j=1}^{m}x_j=0\right\}.
\]
This kernel has the ordered basis
\[
 \alpha_i=e_i-e_{i+1},\qquad 1\leq i\leq m-1,
\]
and the coefficient inverse is explicit:
\begin{equation}
 x=\sum_{i=1}^{m-1}c_i\alpha_i,
 \qquad
 c_i=\sum_{j=1}^{i}x_j.
 \label{eq:global-place-partial-sum-inverse}
\end{equation}
For the standard Euclidean pairing,
\[
 \langle\alpha_i,\alpha_j\rangle
 =\begin{cases}
 2,&i=j,\\
 -1,&|i-j|=1,\\
 0,&|i-j|>1.
 \end{cases}
\]
Thus the balancing kernel, with all coordinates retained, is the root lattice
$A_{m-1}=A_{n+1}$.  More generally, every integer fibre is the affine lattice
\[
 (\delta_P^{\mathrm{bal}})^{-1}(r)
 =-r e_m+\ker\delta_P^{\mathrm{bal}}.
\]
\end{theorem}

\begin{proof}
The formula in \eqref{eq:global-place-balancing-row} shows that the kernel is
the displayed sum-zero lattice.  It is onto because
$\delta_P^{\mathrm{bal}}(-r e_m)=r$ for every $r\in\mathbb Z$.

Each $\alpha_i$ has coordinate sum zero.  Conversely, let
$x=(x_1,\ldots,x_m)$ have sum zero and define the partial sums $c_i$ by
\eqref{eq:global-place-partial-sum-inverse}.  The first coordinate of
$\sum_i c_i\alpha_i$ is $c_1=x_1$; for $2\leq j\leq m-1$ its $j$th coordinate
is
\[
 c_j-c_{j-1}=x_j;
\]
and its last coordinate is
\[
 -c_{m-1}=-\sum_{j=1}^{m-1}x_j=x_m.
\]
This also proves uniqueness of the coefficients, because the first coordinate
recovers $c_1$ and successive coordinate differences recover every $c_i$.
Direct multiplication of the coordinate vectors gives the displayed Cartan
Gram matrix.  Finally, adding $r e_m$ to an element of the fibre over $r$
puts it in the kernel, proving the exact affine-fibre formula.
\end{proof}

For one finite prime, the theorem gives
\[
 \ker\bigl(-(1,1,1):\mathbb Z^3\to\mathbb Z\bigr)
 =\mathbb Z(1,-1,0)\oplus\mathbb Z(0,1,-1)\cong A_2.
\]
This was initially the conjectural stopping point.  The missing codiagonal can,
however, be constructed at the level of stable $C^*$-extensions.  We give the
construction rather than infer it from the displayed integer row.

Put
\[
 \mathcal K=\mathcal K\bigl(\ell^2(\mathbb N_0)\bigr),
 \qquad J=C_0(\mathbb R)\otimes\mathcal K,
\]
and let $e_{00}\in\mathcal K$ be the projection onto the first basis vector.
If
\[
 c(t)=\frac{t-i}{t+i},
\]
then the declared positive generator of $K_1(J)$ is represented in the
unitization by
\begin{equation}
 \gamma(t)=1+(c(t)-1)e_{00}.
 \label{eq:global-place-common-generator}
\end{equation}

\begin{lemma}[A common oriented stabilization]
\label{lem:global-place-common-stabilization}
For every finite prime $p$ and for the archimedean place $\infty$, the proved
local extension admits a stabilized transport
\begin{equation}
 0\longrightarrow J\longrightarrow\widehat E_v
 \longrightarrow\widehat A_v\longrightarrow0,
 \qquad v\in\{p,\infty\},
 \label{eq:global-place-stabilized-local-extension}
\end{equation}
and a morphism from the original local extension to
\eqref{eq:global-place-stabilized-local-extension}.  The quotient algebras are
\[
 \widehat A_p=Q_p\otimes\mathcal K,
 \qquad
 \widehat A_\infty=D_2\otimes\mathcal K.
\]
Under the corner maps on the quotients and the generator
\eqref{eq:global-place-common-generator}, their connecting maps remain
\[
 [1]_p\longmapsto-[\gamma],
 \qquad
 (a[p_+]+b[p_-])\longmapsto-(a+b)[\gamma].
\]
Every stabilization map, generator, and possible orientation reversal is
retained in this assertion.
\end{lemma}

\begin{proof}
For a finite prime, the extension-level Green--Rieffel module in the preceding
subsection gives $I_p\sim_M C_0(\mathbb R)$.  Both algebras are separable and
therefore have countable approximate identities.  The Brown--Green--Rieffel
stabilization theorem gives
\[
 I_p\otimes\mathcal K\cong C_0(\mathbb R)\otimes\mathcal K=J
\]
\cite[Theorems~1.1--1.2, pp.~349--352]{brownGreenRieffel1977}; equivalently,
it follows by taking compact operators after Kasparov stabilization
\cite[Theorem~13.6.2 and Exercise~13.7.1, pp.~131--132]{blackadar1986}.
At the archimedean place, the ideal is
$C_0((0,\infty),M_2)$; the logarithm $(0,\infty)\to\mathbb R$ and a fixed
unitary $\mathbb C^2\otimes\ell^2(\mathbb N_0)\cong
\ell^2(\mathbb N_0)$ give the required stable isomorphism explicitly.

Write $B_v$ for the original local ideal and choose an initial isomorphism
$\theta_v^0:B_v\otimes\mathcal K\to J$.  Its map on the already declared
$K_1$ coordinate is multiplication by a recorded sign
$\varepsilon_v\in\{1,-1\}$.  If $\varepsilon_v=-1$, compose $\theta_v^0$
with the explicit reflection automorphism
\[
 \rho:J\longrightarrow J,\qquad (\rho f)(t)=f(-t),
\]
which takes $[\gamma]$ to $-[\gamma]$; if $\varepsilon_v=1$, leave it
unchanged.  Denote the resulting orientation-preserving isomorphism by
$\theta_v$.  This records the sign correction instead of silently discarding
it.

Tensor the original exact sequence with the nuclear algebra $\mathcal K$.
Let $\tau_v^{\mathrm s}$ be its Busby invariant.  The isomorphism $\theta_v$
extends to multiplier algebras and induces
$\dot\theta_v:Q(B_v\otimes\mathcal K)\to Q(J)$.  Define
\[
 \widehat\tau_v=\dot\theta_v\circ\tau_v^{\mathrm s}
 :\widehat A_v\longrightarrow Q(J)
\]
and let $\widehat E_v$ be its Busby pullback.  The pullback reconstruction and
the induced extension morphism are the constructions of
Sections~15.2--15.3 of Blackadar~\cite[pp.~144--145]{blackadar1986}.  Composing
with the original corner map $x\mapsto x\otimes e_{00}$ gives the asserted
morphism from the unstabilized extension.  Stability sends each displayed
quotient generator to its corresponding rank-one corner, while the declared
$\theta_v$ sends the ideal generator to $[\gamma]$; hence the two local
boundary formulas, including their minus signs, are unchanged.
\end{proof}

Let $\mathcal V=\{\infty\}\cup\{\text{rational primes}\}$ and fix once and for
all a bijection $\nu:\mathcal V\to\mathbb N_0$.  For
\[
 \kappa(j,k)=\frac{(j+k)(j+k+1)}2+k
\]
define an isometry $S_j$ on $\ell^2(\mathbb N_0)$ by
$S_je_k=e_{\kappa(j,k)}$.  Since $\kappa$ is a bijection
$\mathbb N_0^2\to\mathbb N_0$,
\begin{equation}
 S_j^*S_l=\delta_{jl}1,
 \qquad
 \sum_{j=0}^{\infty}S_jS_j^*=1
 \quad\text{strongly}.
 \label{eq:global-place-cuntz-family}
\end{equation}
Write $s_v=1\otimes S_{\nu(v)}\in M(J)$.

\begin{theorem}[Stable place-balancing extension]
\label{thm:global-place-stable-extension}
For every finite prime list $P=(p_1,\ldots,p_n)$ there is an explicitly
constructed extension
\begin{equation}
 0\longrightarrow J\longrightarrow\mathcal E_P
 \longrightarrow\mathcal A_P\longrightarrow0,
 \qquad
 \mathcal A_P=
 \left(\bigoplus_{i=1}^{n}\widehat A_{p_i}\right)
 \oplus\widehat A_\infty,
 \label{eq:global-place-stable-extension}
\end{equation}
with a morphism from the direct sum of the original local extensions after the
declared corner stabilizations.  It has all of the following exact properties.
\begin{enumerate}
\item Its quotient map retains the ordered classes
\[
 ([1]_{p_1},\ldots,[1]_{p_n},[p_+],[p_-])
\]
without merging primes or the two archimedean characters.
\item The ideal map from the direct sum of the stabilized local ideals induces
exactly the codiagonal $\Sigma_P$ of
\eqref{eq:global-place-codiagonal} on $K_1$.
\item The connecting map of \eqref{eq:global-place-stable-extension} is
\[
 \partial_P(k_1,\ldots,k_n,n_+,n_-)
 =-\left(\sum_{i=1}^{n}k_i+n_++n_-\right),
\]
with the displayed sign and all source coordinates retained.  Hence its kernel
is the $A_{n+1}$ lattice of
Theorem~\ref{thm:global-place-balancing-lattice}.
\item Reordering $P$ merely permutes the named quotient summands.  If
$P\subset P'$, zero-inclusion of the new quotient coordinates lifts to a
morphism from the $P$ extension to the $P'$ extension which is the identity on
$J$.  Thus adjoining a prime and removing it by restriction to the zero
coordinate are compatible at extension level, not only on $K$-groups.
\end{enumerate}
\end{theorem}

\begin{proof}
Put $V_P=\{p_1,\ldots,p_n,\infty\}$ and form the finite direct sums
\[
 B_P=\bigoplus_{v\in V_P}J,
 \qquad
 \widehat E_P^\oplus=\bigoplus_{v\in V_P}\widehat E_v,
 \qquad
 \mathcal A_P=\bigoplus_{v\in V_P}\widehat A_v.
\]
For $(b_v)_{v\in V_P}\in B_P$, define
\begin{equation}
 \mu_P((b_v)_v)=\sum_{v\in V_P}s_vb_vs_v^*\in J.
 \label{eq:global-place-ideal-map}
\end{equation}
The orthogonality in \eqref{eq:global-place-cuntz-family} gives, coordinate by
coordinate,
\[
 \mu_P(b)\mu_P(d)
 =\sum_{v,w}s_vb_vs_v^*s_wd_ws_w^*
 =\sum_vs_vb_vd_vs_v^*=\mu_P(bd),
\]
and $s_v^*\mu_P(b)s_v=b_v$.  Thus $\mu_P$ is an injective
$*$-homomorphism.  The same finite formula defines
\[
 \overline\mu_P:M(B_P)=\bigoplus_{v\in V_P}M(J)\longrightarrow M(J).
\]
It maps $B_P$ into $J$ and therefore induces a corona homomorphism
$\dot\mu_P:Q(B_P)\to Q(J)$.

Let
\[
 \widehat\tau_P^\oplus=\bigoplus_{v\in V_P}\widehat\tau_v:
 \mathcal A_P\longrightarrow Q(B_P)
\]
be the Busby invariant of the direct-sum extension and set
\begin{equation}
 \tau_P=\dot\mu_P\circ\widehat\tau_P^\oplus.
 \label{eq:global-place-busby-sum}
\end{equation}
The required middle algebra is the explicit pullback
\begin{equation}
 \mathcal E_P=
 \{(a,m)\in\mathcal A_P\oplus M(J):
       \tau_P(a)=\pi_J(m)\},
 \label{eq:global-place-busby-pullback}
\end{equation}
where $\pi_J:M(J)\to Q(J)$ is the quotient map.  Projection onto $a$ is onto,
and its kernel is exactly $\{(0,j):j\in J\}$, proving
\eqref{eq:global-place-stable-extension}.  In the Busby pullback models for the
local extensions, the middle map is
\begin{equation}
 ((a_v,m_v))_{v\in V_P}
 \longmapsto
 \left((a_v)_{v\in V_P},
       \sum_{v\in V_P}s_vm_vs_v^*\right).
 \label{eq:global-place-middle-map}
\end{equation}
Together with $\mu_P$ on ideals and the identity on $\mathcal A_P$, this is a
morphism of short exact sequences.  It is an extension-level construction
before any $K$-group calculation.  It is the finite orthogonal-corner form of
the standard Busby sum for stable ideals
\cite[Sections~15.6 and 15.9, pp.~149--150, 156]{blackadar1986}.

It remains to compute the ideal map without appealing to its intended answer.
The $v$th copy of \eqref{eq:global-place-common-generator} is sent to
\[
 1+(c(t)-1)q_v,
 \qquad q_v=S_{\nu(v)}e_{00}S_{\nu(v)}^*.
\]
The projection $q_v$ is rank one.  A rotation in the at most two-dimensional
span of $e_0$ and $S_{\nu(v)}e_0$ gives an explicit norm-continuous path of
rank-one projections from $e_{00}$ to $q_v$.  Substituting that path for
$e_{00}$ in \eqref{eq:global-place-common-generator} proves that the $v$th
corner map takes $[\gamma]$ to $[\gamma]$.  Since
$K_1(B_P)=\bigoplus_{v\in V_P}\mathbb Z[\gamma]$, additivity gives
\[
 (\mu_P)_*((y_v)_v)=\left(\sum_{v\in V_P}y_v\right)[\gamma].
\]
This is exactly $\Sigma_P$, not an unidentified isomorphism.

The top extension in the morphism diagram has boundary
$\Delta_P$ from \eqref{eq:global-place-direct-sum-boundary}.  The connecting
maps in $K$-theory are natural for morphisms of short exact sequences
\cite[Definition~21.1.1 and Example~21.1.2(a), pp.~266, 268]{blackadar1986}.
Consequently
\[
 \partial_P=(\mu_P)_*\circ\Delta_P
 =\Sigma_P\circ\Delta_P=\delta_P^{\mathrm{bal}},
\]
which proves the first three assertions and invokes the already proved exact
lattice theorem for the kernel.

Finally, the operators $s_v$ and the transports $\theta_v$ were fixed by the
place $v$, not by its position in $P$.  Reordering therefore only permutes the
direct-sum notation.  If $P\subset P'$, let
$\iota_{P,P'}:\mathcal A_P\to\mathcal A_{P'}$ insert zero in the new place
coordinates.  The literal formulas give
\[
 \tau_{P'}\circ\iota_{P,P'}=\tau_P.
\]
Hence $(a,m)\mapsto(\iota_{P,P'}(a),m)$ maps
$\mathcal E_P$ to $\mathcal E_{P'}$, is the identity on $J$, and makes the
extension diagram commute.  Restriction back to the zero-coordinate summand
recovers the $P$ extension.  This proves the fourth assertion.
\end{proof}

\paragraph{Passage to the natural arithmetic representative.}
The extension-existence, codiagonal, sign, kernel, permutation, and finite-set
compatibility assertions are therefore theorems after stabilization.  The
audited source itself does not define its larger symbol $A_{\varnothing p}$ or
a simultaneous multi-place action.  Subsection
\ref{subsec:unstabilized-semilocal-one-zero} therefore constructs the natural
unit-retaining semilocal crossed product independently and proves its exact
$K$-theory comparison with this stable scalar extension.  Subsection
\ref{subsec:stable-to-natural-extension} then constructs the formerly missing
extension-level $*$-homomorphism by an exact multiplier/corona Busby equation.
The natural boundary still collapses the two archimedean character classes in
$K_0$, and the natural ideal still contains nonconstant unit-fibre classes;
the new theorem records these as the kernel and image data of an injective
stable extension morphism.  No endpoint statement about zeta zeros follows
from any of these extension theorems.

\paragraph{Literature and novelty boundary.}
The proof uses the content-read stable-isomorphism theorem of Brown, Green, and
Rieffel.  It also uses Blackadar's Busby pullback, stable extension sum and
functoriality, and the naturality of $K$-theory connecting maps.  A bounded canonical-index query
for these exact dependencies routed the cited sources.  The earlier bounded
query for the combined terms ``codiagonal'', ``$A_2$'', ``boundary'', and
``$K$-theory'', followed by a phrase-level arXiv check, did not locate this
place-indexed diagram.  That search scope supports no general novelty claim.
The theorem above is the explicit stable construction.  The next two
subsections supply the independently defined arithmetic target and the exact
stable-to-natural extension morphism.  The bounded search scope supports no
general novelty claim.

\paragraph{Formal finite-coordinate replay.}
The file
\begin{center}
\footnotesize
\path{formal/lean/ZetaFunctionFoundation/GlobalPlaceBalancing.lean}
\end{center}
certifies the direct-sum boundary, codiagonal composite, scalar-kernel equation,
surjectivity, and the naturality-implied formula for an arbitrary finite number
of prime coordinates.  Receipt
\path{formal/lean/receipts/FORMAL-20260830-0036.receipt.json} records Lean
4.32.0, Mathlib 4.32.0, exit code zero, and a peak aggregate working set below
the authorized five-gigabyte ceiling.  The Busby invariant, multiplier-algebra
construction, and $C^*$-algebraic naturality theorem remain certified by the
complete standard proof above and its human literature dependencies; this
Lean file does not pretend to formalize those analytic objects.


\subsection{The unstabilized semilocal one-zero algebra: explicit coordinates,
boundary maps, and information loss}
\label{subsec:unstabilized-semilocal-one-zero}

The audited source introduces $A_{\varnothing p}$ only as an
``intermediate semilocal crossed-product algebra.''  It does not specify its
space, action, completion, inclusion, quotient, or imprimitivity module.  The
two-place construction of Connes, Consani, and Marcolli does, however, supply
the exact seed: $\mathbb Q_p\times\mathbb R$, acted on diagonally by
$\{\pm p^n:n\in\mathbb Z\}$, with generic and one-zero strata
\cite[Section~8.2]{connesConsaniMarcolli2007}.  We now perform the finite-place
construction directly.  This gives a natural meaning to the missing symbol;
it is not an assertion that this meaning was already present in the local
source or is the unique possible enrichment of that symbol.

Let $P$ be a nonempty finite set of rational primes and put
\[
 S_P=P\sqcup\{\infty\},\qquad
 \mathbb A_{S_P}=\left(\prod_{q\in P}\mathbb Q_q\right)\times\mathbb R,
\]
\[
 \Gamma_P=\mathbb Z[1/\!\prod_{q\in P}q]^{\times}
 =\left\{\varepsilon\prod_{q\in P}q^{n_q}:
   \varepsilon\in\{\pm1\},\ n_q\in\mathbb Z\right\}.
\]
The group $\Gamma_P\cong(\mathbb Z/2)\times\mathbb Z^P$ acts by diagonal
multiplication.  For $x=((x_q)_{q\in P},y)\in\mathbb A_{S_P}$ define its
literal zero set
\[
 Z(x)=\{q\in P:x_q=0\}\cup
 \begin{cases}\{\infty\},&y=0,\\ \varnothing,&y\ne0.\end{cases}
\]
No support coordinate is identified with another one in this definition.
Set
\[
 W_P^{(1)}=\{x:\#Z(x)\leq1\},\qquad
 W_{\varnothing}=\{x:Z(x)=\varnothing\},\qquad
 W_v=\{x:Z(x)=\{v\}\}.
\]
The complement of $W_P^{(1)}$ is a finite union of closed coordinate
intersections, so $W_P^{(1)}$ is an invariant locally compact open subspace.

\begin{definition}[Natural unstabilized one-zero algebra]
\label{def:unstabilized-one-zero-algebra}
All crossed products in this subsection are full crossed products.  Define
\begin{align*}
 \mathfrak A_P^{(1)}&=C_0(W_P^{(1)})\rtimes\Gamma_P,\\
 \mathfrak I_P&=C_0(W_{\varnothing})\rtimes\Gamma_P,\\
 \mathfrak D_{v,P}&=C_0(W_v)\rtimes\Gamma_P,\qquad
 \mathfrak D_P=\bigoplus_{v\in S_P}\mathfrak D_{v,P}.
\end{align*}
For $p\in P$, the source symbol receives the explicit candidate
\[
 A_{\varnothing p}^{(P)}
 :=C_0(W_{\varnothing}\cup W_p)\rtimes\Gamma_P,
\]
with $A_{\varnothing}^{(P)}=\mathfrak I_P$ and
$A_p^{(P)}=\mathfrak D_{p,P}$.
\end{definition}

\begin{theorem}[The literal one-zero extension]
\label{thm:unstabilized-one-zero-extension}
Restriction to the closed one-zero boundary gives the short exact sequence
\begin{equation}
 0\longrightarrow\mathfrak I_P
 \longrightarrow\mathfrak A_P^{(1)}
 \longrightarrow\mathfrak D_P\longrightarrow0.
 \label{eq:unstabilized-one-zero-extension}
\end{equation}
For every $v\in S_P$, the invariant open subspace
$W_{\varnothing}\cup W_v$ gives a morphism from the $v$th edge extension
to \eqref{eq:unstabilized-one-zero-extension}, equal to the identity on
$\mathfrak I_P$ and to the summand inclusion on the quotient.
Full and reduced versions of all these algebras agree.
\end{theorem}

\begin{proof}
Inside $W_P^{(1)}$, the generic stratum is open and its complement is the
finite disjoint union $\bigsqcup_{v\in S_P}W_v$.  Restriction of functions
therefore gives
\[
 0\to C_0(W_{\varnothing})\to C_0(W_P^{(1)})
 \to\bigoplus_{v\in S_P}C_0(W_v)\to0.
\]
The maps are $\Gamma_P$-equivariant, and the full crossed-product functor is
exact.  This proves \eqref{eq:unstabilized-one-zero-extension}.  Removing all
boundary strata except $W_v$ leaves an invariant open subspace; extension by
zero gives the asserted middle and quotient maps and is literally the
identity on the common generic ideal.  Finally, $\Gamma_P$ is abelian-by-finite
and hence amenable.  Exel's approximation and regular-representation theorems
identify the full and reduced cross-sectional algebras in this case
\cite[Theorems~20.6--20.7]{exel2017}.
\end{proof}

We next retain every unit coordinate and calculate a finite-prime edge.  Put
\[
 K_P=\prod_{q\in P}\mathbb Z_q^{\times},\qquad
 \ell_q=\log q.
\]
For $p\in P$, let $K_{P\setminus p}=\prod_{q\in P\setminus\{p\}}
\mathbb Z_q^{\times}$, with the empty product understood as one point, and
define
\[
 \beta_p((b_q)_{q\ne p},s)
 =((p b_q)_{q\ne p},s+\ell_p).
\]
The finite-prime boundary carrier is the compact mapping torus
\begin{equation}
 B_{p,P}=(K_{P\setminus p}\times\mathbb R)/\langle\beta_p\rangle.
 \label{eq:multiplace-finite-mapping-torus}
\end{equation}

\begin{theorem}[Complete coordinates for the $P$-semilocal $p$-edge]
\label{thm:multiplace-finite-edge-coordinates}
The action of $\Gamma_P$ on $W_{\varnothing}\cup W_p$ is free and proper.
As a set, its orbit space is
\begin{equation}
 \mathcal S_{P,p}=(K_P\times\mathbb R)\sqcup B_{p,P}.
 \label{eq:multiplace-finite-edge-space}
\end{equation}
The symbol $\sqcup$ here records the two strata, not the topological-sum
topology; the cross-stratum convergence is specified below.
For a generic point, write
\[
 m_q=v_q(x_q),\qquad u_q=q^{-m_q}x_q\in\mathbb Z_q^{\times},
 \qquad\epsilon=\sgn(y).
\]
Its orbit coordinates are, without suppressing any cross-prime factor,
\begin{equation}
 a_q=\epsilon\!\prod_{r\in P\setminus\{q\}}r^{-m_r}u_q,
 \qquad
 t=\log|y|-\sum_{r\in P}m_r\ell_r.
 \label{eq:multiplace-generic-coordinates}
\end{equation}
For a point of $W_p$, define for $q\ne p$
\begin{equation}
 b_q=\epsilon\!\prod_{r\in P\setminus\{p,q\}}r^{-m_r}u_q,
 \qquad
 s=\log|y|-\sum_{r\ne p}m_r\ell_r;
 \label{eq:multiplace-p-boundary-coordinates}
\end{equation}
its orbit coordinate is $[(b_q)_{q\ne p},s]\in B_{p,P}$.

A generic net $(a^{(j)},t_j)$ converges to $[b,s]\in B_{p,P}$ exactly when
\begin{equation}
 t_j\longrightarrow-\infty,\qquad
 [(a_q^{(j)})_{q\ne p},t_j]\longrightarrow[b,s]
 \quad\hbox{in }B_{p,P};
 \label{eq:multiplace-finite-cross-convergence}
\end{equation}
there is no condition on the discarded coordinate
$a_p^{(j)}\in\mathbb Z_p^{\times}$.
\end{theorem}

\begin{proof}
Let $\gamma=\delta\prod_{r\in P}r^{n_r}\in\Gamma_P$.  On the generic
stratum,
\[
 m_q'=m_q+n_q,\qquad
 u_q'=\delta\!\prod_{r\ne q}r^{n_r}u_q,
 \qquad\epsilon'=\delta\epsilon.
\]
Substitution into \eqref{eq:multiplace-generic-coordinates} cancels every
$n_r$ and both factors $\delta$; the logarithmic coordinate is invariant
because $\log|\gamma|=\sum_rn_r\ell_r$.  Conversely, multiplication by the
unique element
\[
 \epsilon\prod_{r\in P}r^{-m_r}
\]
makes $y$ positive and every finite valuation zero.  Hence $(a,t)$ is a
complete generic orbit coordinate, not only an invariant.

On $W_p$, the same normalization is available for the sign and for all
valuations $m_q$ with $q\ne p$.  The unused power $p^k$ acts on the resulting
coordinates by $\beta_p^k$, which proves the boundary formula and its
uniqueness modulo \eqref{eq:multiplace-finite-mapping-torus}.

For the topology, let a generic representative have $p$-valuation $m_p$.
Equations \eqref{eq:multiplace-generic-coordinates} and
\eqref{eq:multiplace-p-boundary-coordinates} give the literal relations
\begin{equation}
 b_q=p^{m_p}a_q\quad(q\ne p),\qquad s=t+m_p\ell_p.
 \label{eq:multiplace-edge-coordinate-change}
\end{equation}
Thus $[(b_q),s]=[(a_q)_{q\ne p},t]$ in $B_{p,P}$.  Convergence of the
$p$-adic coordinate to zero is precisely $m_p\to+\infty$; if a boundary
representative remains in a compact set, the second identity forces
$t\to-\infty$ and gives \eqref{eq:multiplace-finite-cross-convergence}.
Conversely, use a local section of the proper $\mathbb Z$-quotient defining
$B_{p,P}$ to choose integers $m_j$ such that
\[
 \beta_p^{m_j}((a_q^{(j)})_{q\ne p},t_j)\longrightarrow(b,s).
\]
Because $t_j\to-\infty$ while the lifted real coordinates converge to $s$,
necessarily $m_j\to+\infty$.
Take all other valuations zero, set
$u_q^{(j)}=p^{m_j}a_q^{(j)}$ for $q\ne p$,
$x_p^{(j)}=p^{m_j}a_p^{(j)}$, and
$y_j=\exp(t_j+m_j\ell_p)$.  These are valid unit coordinates and converge to
the required boundary representative.  This proves both directions and the
absence of a condition on $a_p^{(j)}$.

The action is free because $y\ne0$.  For properness, let $C,D$ be compact
subsets of the edge and suppose $\gamma C\cap D\ne\varnothing$.  Every
$q\ne p$ coordinate is nonzero throughout the edge, so its valuations on
$C$ and $D$ bound $n_q$.  The real coordinate is also nonzero, so its absolute
values on $C,D$ bound $\sum_qn_q\ell_q$.  The already bounded $n_q$ for
$q\ne p$ then bound $n_p$, and the sign has only two values.  Only finitely
many $\gamma$ occur, proving properness.
\end{proof}

There is now an exact morphism to the one-ended local model, rather than a
Morita identification that erases the unit carrier.  Define
\[
 \rho_{p,P}:B_{p,P}\to S^1_{\ell_p},\qquad
 \rho_{p,P}([(b_q),s])=[s]_{\ell_p},
\]
and
\begin{equation}
 \Pi_{P,p}:\mathcal S_{P,p}\to X_{\ell_p},\qquad
 \Pi_{P,p}(a,t)=t,\quad
 \Pi_{P,p}([b,s])=[s]_{\ell_p}.
 \label{eq:multiplace-to-local-quotient}
\end{equation}
The mapping-torus relation changes $s$ by an integral multiple of $\ell_p$,
so $\rho_{p,P}$ is well defined.  The compact group $K_P$ acts by
\[
 k\cdot(a,t)=(ka,t),\qquad
 k\cdot[b,s]=[(k_qb_q)_{q\ne p},s].
\]
The second formula is well defined because coordinatewise multiplication by
$k$ commutes with $\beta_p$; the factor $k_p$ acts trivially on the boundary.
The convergence criterion proves that this action is continuous and that
$\Pi_{P,p}$ is exactly its orbit map.  Hence $\Pi_{P,p}$ is closed and proper.
Its generic fibres are $K_P$; its boundary fibres are copies of
$K_{P\setminus p}$.  Pullback therefore gives
\begin{equation}
\begin{array}{ccccccccc}
0&\to&C_0(\mathbb R)&\to&C_0(X_{\ell_p})&\to&C(S^1_{\ell_p})&\to&0\\
&&\downarrow\scriptstyle{\pi_P^*}&&\downarrow\scriptstyle{\Pi_{P,p}^*}
&&\downarrow\scriptstyle{\rho_{p,P}^*}\\
0&\to&C_0(K_P\times\mathbb R)&\to&C_0(\mathcal S_{P,p})
&\to&C(B_{p,P})&\to&0,
\end{array}
\label{eq:multiplace-finite-pullback-diagram}
\end{equation}
where $\pi_P(a,t)=t$.  Green imprimitivity, applied with the restricted
modules exactly as in the two-place proof, identifies the lower row with the
crossed-product edge extension
\[
 0\to\mathfrak I_P\to A_{\varnothing p}^{(P)}
 \to\mathfrak D_{p,P}\to0
\]
at the Morita level
\cite[Theorem~1.3.22]{meslandSengun2025}.

\begin{corollary}[Finite-place unit class with every unit fibre retained]
\label{cor:multiplace-finite-boundary-class}
Under
\[
 K_1(\mathfrak I_P)\cong
 K_1(C_0(K_P\times\mathbb R))
 \cong K_0(C(K_P))\cong C(K_P,\mathbb Z),
\]
the connecting map for the $p$-edge sends the Green image
$\eta_{p,P}$ of $[1_{B_{p,P}}]$ to
\begin{equation}
 \partial_{p,P}(\eta_{p,P})=-\mathbf1_{K_P}.
 \label{eq:multiplace-finite-boundary-class}
\end{equation}
\end{corollary}

\begin{proof}
The right vertical arrow of
\eqref{eq:multiplace-finite-pullback-diagram} is unital, so it sends the local
circle unit to the mapping-torus unit.  Theorem~\ref{thm:primitive-local-boundary}
gives $-1$ for the top boundary map.  Naturality sends it through
$\pi_P^*$, which is the constant-function inclusion
$\mathbb Z\to C(K_P,\mathbb Z)$.  This proves
\eqref{eq:multiplace-finite-boundary-class}.  The last $K_0$ identification
retains all locally constant integer-valued rank functions on the profinite
space $K_P$; it is not a normalization to one integer.
\end{proof}

The archimedean edge contains a different, equally explicit, information-loss
mechanism.  Normalize all finite valuations but do not choose a real sign.  If
$x_q=q^{m_q}u_q$, put
\begin{equation}
 c_q=\prod_{r\in P\setminus\{q\}}r^{-m_r}u_q,
 \qquad z=y\prod_{r\in P}r^{-m_r}.
 \label{eq:multiplace-arch-normalization}
\end{equation}
The residual sign acts by $(c,z)\mapsto(-c,-z)$.

\begin{theorem}[Archimedean edge and collapse of the two character labels]
\label{thm:multiplace-arch-boundary}
The action of $\Gamma_P$ on $W_{\varnothing}\cup W_\infty$ is free and
proper.  After Morita reduction of the free valuation directions, its
extension is
\begin{equation}
\begin{array}{ccccccccc}
0&\to&C_0(K_P\times\mathbb R^\times)\rtimes C_2
&\to&C_0(K_P\times\mathbb R)\rtimes C_2
&\to&C(K_P)\rtimes C_2&\to&0,
\end{array}
\label{eq:multiplace-arch-crossed-extension}
\end{equation}
where $C_2$ acts diagonally by $(c,z)\mapsto(-c,-z)$ and by
$c\mapsto-c$ on the boundary.  Constant functions in $K_P$ give a morphism
from the local reflection extension
\eqref{eq:primitive-archimedean-extension} to
\eqref{eq:multiplace-arch-crossed-extension}.  If
\[
 e_\pm=\frac{1\pm U}{2}\in C(K_P)\rtimes C_2,
\]
then
\begin{equation}
 \partial_{\infty,P}[e_+]
 =\partial_{\infty,P}[e_-]
 =-\mathbf1_{K_P}\in C(K_P,\mathbb Z).
 \label{eq:multiplace-arch-boundary-classes}
\end{equation}
Moreover $[e_+]=[e_-]$ in $K_0(C(K_P)\rtimes C_2)$.
\end{theorem}

\begin{proof}
Formula \eqref{eq:multiplace-arch-normalization} is invariant under every
prime-power direction, and every orbit has a unique representative with all
finite valuations zero, up to the residual sign.  This gives
\eqref{eq:multiplace-arch-crossed-extension}.  Freeness follows because all
finite coordinates are nonzero on this edge.  Properness follows because
their valuations on two compact sets bound every prime exponent; the residual
sign is finite.

The projection $K_P\times\mathbb R\to\mathbb R$ is equivariant for the
diagonal and reflection actions.  Pullback, followed by crossed product, is
therefore the promised extension morphism.  It sends the local projections
$p_\pm$ to $e_\pm$ and sends the local ideal generator to the constant rank
function.  Equation \eqref{eq:multiplace-arch-boundary-classes} follows from
Theorem~\ref{thm:primitive-archimedean-boundary} by naturality.

It remains to prove, rather than assume, the equality of the two boundary
classes.  Choose one $q_0\in P$.  If $q_0$ is odd, choose one representative
from every pair $\{r,-r\}$ in $(\mathbb Z/q_0\mathbb Z)^\times$.  If $q_0=2$,
choose the residue class $1\pmod 4$; its negative is the class
$3\pmod 4$.  The inverse image of the chosen classes defines a clopen set
$H\subset K_P$ with
\[
 K_P=H\sqcup(-H).
\]
Thus the principal double cover $K_P\to K_P/\{\pm1\}$ has the explicit
continuous section whose image is $H$.  Consequently
\[
 C(K_P)\rtimes C_2\cong M_2(C(H)),
\]
and $e_+$ and $e_-$ are rank-one projections in the same $2\times2$ matrix
algebra.  In the standard two-sheet matrix coordinates the constant unitary
$\operatorname{diag}(1,-1)$ conjugates $e_+$ to $e_-$, proving
$[e_+]=[e_-]$.  This is the exact point at which the natural semilocal carrier
forgets the local sign-character difference.
\end{proof}

We can now compute the natural multi-boundary map without pretending that the
whole action is proper.  Let
\[
 \Lambda_P=\bigoplus_{p\in P}\mathbb Z\epsilon_p
 \oplus\mathbb Z\epsilon_+\oplus\mathbb Z\epsilon_-.
\]
Define the boundary-label morphism
\begin{equation}
 R_P:\Lambda_P\longrightarrow K_0(\mathfrak D_P),\qquad
 R_P(\epsilon_p)=\eta_{p,P},\quad
 R_P(\epsilon_\pm)=[e_\pm],
 \label{eq:unstabilized-boundary-label-morphism}
\end{equation}
where the last two classes are transported through the valuation Morita
module.  Also define the exact constant-coordinate morphism
\[
 c_P:\mathbb Z\longrightarrow C(K_P,\mathbb Z),\qquad
 c_P(n)=n\mathbf1_{K_P}.
\]

\begin{theorem}[Unstabilized arithmetic realization of the codiagonal shadow]
\label{thm:unstabilized-codiagonal-shadow}
For the connecting map of
\eqref{eq:unstabilized-one-zero-extension},
\begin{equation}
 \partial_P^{(1)}R_P
 ( (k_p)_{p\in P},n_+,n_- )
 =-\left(\sum_{p\in P}k_p+n_++n_-\right)\mathbf1_{K_P}.
 \label{eq:unstabilized-codiagonal-formula}
\end{equation}
Equivalently, the square
\begin{equation}
\begin{CD}
 \Lambda_P @>{\delta_P^{\mathrm{bal}}}>> \mathbb Z\\
 @V{R_P}VV @VV{c_P}V\\
 K_0(\mathfrak D_P) @>{\partial_P^{(1)}}>> K_1(\mathfrak I_P)
\end{CD}
\label{eq:stable-to-unstabilized-k-square}
\end{equation}
commutes under the displayed Morita and suspension coordinates.

The morphism $R_P$ does not preserve all formal labels injectively:
\begin{equation}
 \ker R_P=\mathbb Z(\epsilon_+-\epsilon_-).
 \label{eq:unstabilized-boundary-label-kernel}
\end{equation}
Hence the formal scalar kernel $A_{|P|+1}$ maps onto the actual kernel on the
distinguished boundary-unit span, with exact kernel
$\mathbb Z(\epsilon_+-\epsilon_-)$.  The latter actual kernel has rank $|P|$
and an explicit $A_{|P|}$ basis after replacing the two archimedean labels by
their common image.
\end{theorem}

\begin{proof}
For each $v$, the edge algebra is an invariant open ideal in
$\mathfrak A_P^{(1)}$, and Theorem~\ref{thm:unstabilized-one-zero-extension}
gives a morphism of its extension into the global one-zero extension.  The
connecting maps are natural
\cite[Definition~21.1.1 and Example~21.1.2(a)]{blackadar1986}.  Therefore
Corollary~\ref{cor:multiplace-finite-boundary-class} supplies the term
$-k_p\mathbf1_{K_P}$ for every finite $p$, while
Theorem~\ref{thm:multiplace-arch-boundary} supplies
$-n_+\mathbf1_{K_P}$ and $-n_-\mathbf1_{K_P}$.  Additivity proves
\eqref{eq:unstabilized-codiagonal-formula} and the square.

The quotient $\mathfrak D_P$ is a direct sum over the distinct boundary
strata.  Each finite unit class has nonzero rank in its own summand and is
therefore independent of all other finite and archimedean summands.  In the
archimedean summand, the preceding theorem proves that the two classes agree;
their common rank is one, so this is their only relation.  This proves
\eqref{eq:unstabilized-boundary-label-kernel}.  The formal kernel of
$\delta_P^{\mathrm{bal}}$ is the already proved root lattice
$A_{|P|+1}$.  Quotienting by the displayed simple difference identifies the
two archimedean coordinates and leaves
\[
 \left\{((k_p)_{p\in P},n)\in\mathbb Z^{P}\oplus\mathbb Z:
       \sum_pk_p+n=0\right\},
\]
whose consecutive-difference basis is $A_{|P|}$.  No lattice coordinate has
been silently deleted: the lost coordinate and the map that loses it are
both displayed.
\end{proof}

There is a final analytic distinction.  Although every individual edge action
is free and proper, the simultaneous action on $W_P^{(1)}$ is not proper.

\begin{proposition}[Explicit obstruction to a global Green reduction]
\label{prop:unstabilized-global-action-nonproper}
For every nonempty $P$, the action of $\Gamma_P$ on $W_P^{(1)}$ is not proper.
Consequently one may not replace
\eqref{eq:unstabilized-one-zero-extension} by the commutative orbit-space
extension using Green imprimitivity.  This does not affect its crossed-product
exactness or the edgewise boundary calculation above.
\end{proposition}

\begin{proof}
Fix $p\in P$.  Choose integers $n_j\to\infty$ along a subsequence for which
the tuple $(p^{n_j})_{q\in P\setminus\{p\}}$ converges in the compact group
$K_{P\setminus p}$.  Put
\[
 z_j=((x_{q,j})_{q\in P},y_j),\qquad
 x_{p,j}=p^{n_j},\quad x_{q,j}=p^{n_j}\ (q\ne p),\quad y_j=1.
\]
Then $z_j$ converges to a point of the $p$-boundary.  For
$\gamma_j=p^{-n_j}$,
\[
 \gamma_jz_j=((1)_{q\in P},p^{-n_j}),
\]
which converges to a point of the archimedean boundary.  The union of both
sequences and their limits is compact in $W_P^{(1)}$, while the distinct
$\gamma_j$ escape every finite subset of the discrete group $\Gamma_P$.
Thus the inverse image of a compact subset of
$W_P^{(1)}\times W_P^{(1)}$ under $(\gamma,z)\mapsto(\gamma z,z)$ is not
compact.  This is exactly failure of properness.
\end{proof}

\paragraph{Exact comparison boundary before the extension lift.}
The natural algebra \eqref{eq:unstabilized-one-zero-extension} is now an
independently defined unstabilized arithmetic representative, and
\eqref{eq:stable-to-unstabilized-k-square} is an exact morphism from the
stable scalar boundary calculation to its unit-retaining $K$-theory shadow.
It is not an isomorphism: the explicit morphism
$R_P$ has kernel \eqref{eq:unstabilized-boundary-label-kernel}, while the ideal
map $c_P$ retains only the constant subgroup of the much larger group
$C(K_P,\mathbb Z)$.  This proves the precise information loss; it does not
assert that the two constructions are categorically unrelated.

Subsection \ref{subsec:stable-to-natural-extension} resolves that narrower
problem after stabilization.  It writes the ideal, quotient, middle,
multiplier, and corona maps explicitly, proves the Busby equation, and gives a
minimal symmetric split enrichment when the two sign classes must remain
distinct in $K_0$.  The nonconstant unit-fibre classes remain outside the
stable scalar image and are retained as such.  No statement about an
individual zeta zero follows from the construction.

\paragraph{Bounded coordinate replay.}
The script
\begin{center}
\footnotesize
\path{certificates/globalization_nte_multiplace_unstabilized/}\\
\path{verify_multiplace_unstabilized.py}
\end{center}
replays 362 named checks for one through eight finite primes.  They cover
exponents and signs, screw coordinates, the lattice square, the clopen
section, and the nonproperness sequence.  The accompanying JSON receipt has
stable ID
\begin{center}
\footnotesize
\texttt{CERT-GNTE-MULTIPLACE-}\\
\texttt{UNSTABILIZED-20260830-0001}.
\end{center}
It records two Python~3.13/SymPy runs with exit code zero.  It does not replace
the printed topology, crossed-product, Green-imprimitivity, or six-term
naturality proofs, and it launches no Lean process.


\subsection{Lifting the balancing square: an exact stable embedding into the
natural one-zero extension}
\label{subsec:stable-to-natural-extension}

The preceding subsection left an extension-level comparison as an open
problem.  The comparison does exist after stabilizing the natural one-zero
extension.  The compact unit fibre and the archimedean sign collapse do not
obstruct the morphism.  Instead, they become its exact image and kernel data.
We construct the ideal, quotient, middle, multiplier, and corona maps before
using their induced maps on $K$-theory.

Retain a nonempty finite prime set $P$.  Abbreviate
\[
 \mathfrak N_P=\mathfrak A_P^{(1)},\qquad
 \mathfrak I_P^{\mathrm s}=\mathfrak I_P\otimes\mathcal K,\qquad
 \mathfrak N_P^{\mathrm s}=\mathfrak N_P\otimes\mathcal K,\qquad
 \mathfrak D_P^{\mathrm s}=\mathfrak D_P\otimes\mathcal K.
\]
Tensoring \eqref{eq:unstabilized-one-zero-extension} with the nuclear algebra
$\mathcal K$ gives
\begin{equation}
 0\longrightarrow\mathfrak I_P^{\mathrm s}
 \longrightarrow\mathfrak N_P^{\mathrm s}
 \xrightarrow{q_P^{\mathrm s}}\mathfrak D_P^{\mathrm s}
 \longrightarrow0.
 \label{eq:natural-one-zero-stabilized-extension}
\end{equation}

\begin{lemma}[Essential generic ideal and its literal Busby pullback]
\label{lem:natural-one-zero-essential-pullback}
The ideal $\mathfrak I_P$ is essential in $\mathfrak N_P$.  Consequently
$\mathfrak I_P^{\mathrm s}$ is essential in
$\mathfrak N_P^{\mathrm s}$, and the latter is canonically isomorphic to
\begin{equation}
 \left\{(d,n)\in\mathfrak D_P^{\mathrm s}
       \oplus M(\mathfrak I_P^{\mathrm s}):
       \tau_P^{\mathrm{nat}}(d)=\pi_{\mathfrak I}(n)\right\},
 \label{eq:natural-one-zero-busby-pullback}
\end{equation}
where
$\tau_P^{\mathrm{nat}}:\mathfrak D_P^{\mathrm s}
\to Q(\mathfrak I_P^{\mathrm s})$ is the Busby invariant and
$\pi_{\mathfrak I}$ is the multiplier quotient.
\end{lemma}

\begin{proof}
Every one-zero boundary point is a limit of points for which the vanishing
coordinate has been replaced by a nonzero coordinate tending to zero.  Thus
$W_{\varnothing}$ is dense in $W_P^{(1)}$, and
$C_0(W_{\varnothing})$ is an essential ideal of
$C_0(W_P^{(1)})$.

It remains to check that crossing with $\Gamma_P$ does not lose essentiality.
Use the reduced crossed product, which equals the full crossed product because
$\Gamma_P$ is amenable.  Let $E$ be its faithful conditional expectation and,
for $g\in\Gamma_P$, write
\[
 x_g=E(xu_g^*)\in C_0(W_P^{(1)})
\]
for the $g$th Fourier coefficient.  Suppose that
$x\mathfrak I_P=0$.  For every
$f\in C_0(W_{\varnothing})$,
\[
 0=E(xfu_g^*)=x_g\,\alpha_g(f).
\]
The invariant ideal $C_0(W_{\varnothing})$ is essential, so $x_g=0$ for
every $g$.  Faithfulness of the regular Fourier representation gives $x=0$.
Applying the same argument to $x^*$ removes the right annihilator as well.
Hence $\mathfrak I_P$ is essential.

Stabilization preserves essentiality.  The canonical homomorphism
$\mathfrak N_P^{\mathrm s}\to M(\mathfrak I_P^{\mathrm s})$ is therefore
injective.  Busby reconstruction identifies an essential extension with the
pullback of its Busby invariant and the multiplier quotient
\cite[Sections~15.2--15.3, pp.~144--145]{blackadar1986}, proving
\eqref{eq:natural-one-zero-busby-pullback}.
\end{proof}

The generic action is free and proper, and the complete orbit coordinates of
Theorem~\ref{thm:multiplace-finite-edge-coordinates} identify its orbit space
with $K_P\times\mathbb R$.  The Green module therefore gives
\[
 \mathfrak I_P\sim_M C_0(K_P\times\mathbb R)
\]
\cite[Theorem~1.3.22]{meslandSengun2025}.  Fix one stabilization isomorphism,
constructed from that module and its linking algebra,
\begin{equation}
 \chi_P:\mathfrak I_P^{\mathrm s}
 \xrightarrow{\ \cong\ }
 C(K_P)\otimes C_0(\mathbb R)\otimes\mathcal K
 =C(K_P)\otimes J.
 \label{eq:natural-generic-stable-coordinate}
\end{equation}
This is a declared coordinate choice.  It does not identify $K_P$ with a
point.

\begin{definition}[Constant-unit-fibre ideal embedding]
\label{def:constant-unit-fibre-ideal-embedding}
Define the nondegenerate injective $*$-homomorphism
\begin{equation}
 \jmath_P:J\longrightarrow\mathfrak I_P^{\mathrm s},
 \qquad
 \chi_P(\jmath_P(b))=1_{C(K_P)}\otimes b.
 \label{eq:constant-unit-fibre-ideal-embedding}
\end{equation}
Write
$\overline\jmath_P:M(J)\to M(\mathfrak I_P^{\mathrm s})$ and
$\dot\jmath_P:Q(J)\to Q(\mathfrak I_P^{\mathrm s})$ for its multiplier
and corona extensions.
\end{definition}

For every $k\in K_P$, evaluation at $k$ after $\chi_P$ gives a
$*$-homomorphism
\[
 \epsilon_{P,k}:\mathfrak I_P^{\mathrm s}\longrightarrow J,
 \qquad \epsilon_{P,k}\jmath_P=\operatorname{id}_J.
\]
It is nondegenerate and extends to multipliers and coronas.  Hence
$\jmath_P$, $\overline\jmath_P$, and $\dot\jmath_P$ are all split
injective.  On the declared $K_1$ coordinates,
\begin{equation}
 (\jmath_P)_*(n[\gamma])
 =n\mathbf1_{K_P},
 \label{eq:constant-unit-fibre-K1-map}
\end{equation}
which is exactly the map $c_P$ in
\eqref{eq:stable-to-unstabilized-k-square}.

We next make the placewise edge maps compatible before adding them.

\begin{lemma}[Compatible stabilized edge morphisms]
\label{lem:compatible-stabilized-edge-morphisms}
For each $v\in V_P=P\sqcup\{\infty\}$, the stable transports in
Lemma~\ref{lem:global-place-common-stabilization} may be chosen together with
injective $*$-homomorphisms
\[
 \lambda_{v,P}:\widehat E_v\longrightarrow\mathfrak N_P^{\mathrm s},
 \qquad
 \alpha_{v,P}:\widehat A_v\longrightarrow
 \mathfrak D_{v,P}\otimes\mathcal K
 \subset\mathfrak D_P^{\mathrm s},
\]
so that
\begin{equation}
\begin{CD}
0@>>>J@>>>\widehat E_v@>>>\widehat A_v@>>>0\\
@.@V{\jmath_P}VV@V{\lambda_{v,P}}VV@V{\alpha_{v,P}}VV@.\\
0@>>>\mathfrak I_P^{\mathrm s}
 @>>>\mathfrak N_P^{\mathrm s}
 @>{q_P^{\mathrm s}}>>\mathfrak D_P^{\mathrm s}@>>>0
\end{CD}
\label{eq:compatible-stabilized-edge-morphism}
\end{equation}
commutes.  Equivalently, its Busby equation is
\begin{equation}
 \tau_P^{\mathrm{nat}}\alpha_{v,P}
 =\dot\jmath_P\widehat\tau_v.
 \label{eq:placewise-natural-busby-equation}
\end{equation}
At a finite place, the quotient map is the stabilized form of
$\rho_{p,P}^*$ in
\eqref{eq:multiplace-finite-pullback-diagram}.  At infinity it sends
$p_\pm$ to $e_\pm$ as in
Theorem~\ref{thm:multiplace-arch-boundary}.
\end{lemma}

\begin{proof}
We first expose the coordinate which forces the ideal arrow.  The generic
normalization in
\eqref{eq:multiplace-generic-coordinates} is a global section of the
principal $\Gamma_P$-bundle: the representative of $(a,t)$ is
\[
 ((a_q)_{q\in P},e^t)\in
 \left(\prod_{q\in P}\mathbb Q_q^\times\right)\times\mathbb R^\times.
\]
Consequently the generic orbit is literally $K_P\times\mathbb R$.  Every
finite local-to-semilocal orbit map restricts on that stratum to
\[
 \pi_P:K_P\times\mathbb R\longrightarrow\mathbb R,
 \qquad (a,t)\longmapsto t,
\]
and its function pullback is exactly
$f(t)\mapsto\mathbf1_{K_P}(a)f(t)$.  At infinity, the quotient of
$K_P\times\mathbb R^\times$ by $(c,z)\mapsto(-c,-z)$ has the equally
explicit section $z>0$ and coordinates $(a,t)=(\operatorname{sgn}(z)c,
\log|z|)$.  The quotient of the equivariant projection to
$\mathbb R^\times$ is again $(a,t)\mapsto t$.  Thus all placewise ideal
arrows are the same constant-unit-fibre map before stabilization; this is not
an appeal to equality of their induced $K$-classes.

For finite $p$, the two commutative rows in
\eqref{eq:multiplace-finite-pullback-diagram} form a literal injective
morphism of extensions.  The local crossed-product row and the upper
commutative row are Morita equivalent at extension level by the restricted
Green modules used in the proof of
Theorem~\ref{thm:primitive-local-boundary}; the semilocal edge row and the
lower commutative row are Morita equivalent by the same construction.  At
infinity, the equivariant projection
$K_P\times\mathbb R\to\mathbb R$ gives the literal crossed-product morphism,
and its generic free-action Green module has the orbit map just calculated.

We now stabilize these \emph{extension triples}, rather than three unrelated
algebras.  The linking algebra of each extension-level Green module contains
the linking algebra of the restricted ideal as a closed ideal and has the
quotient-module linking algebra as its quotient.  After tensoring with
$\mathcal K$, Brown--Green--Rieffel equivalence of the two complementary full
linking projections supplies one multiplier partial isometry.  Conjugation by
that partial isometry gives stable isomorphisms of the ideal, middle, and
quotient corners simultaneously.  Since multipliers carry the linking ideal
into itself, the same conjugation restricts to the ideal and descends to the
quotient, so the three stable isomorphisms form a commuting extension
diagram.  This is the linking-algebra construction in
\cite[Theorems~1.1--1.2, pp.~349--352]{brownGreenRieffel1977}; the compact
operator model of the Green modules is
\cite[Theorem~1.3.22, p.~14]{meslandSengun2025}.

Choose the source transports in
Lemma~\ref{lem:global-place-common-stabilization} and the target coordinate
$\chi_P$ by these simultaneous stable diagrams, using fixed Hilbert-space
unitaries for their $\mathcal K$ factors.  In those declared coordinates the
ideal component of the composite is the already calculated map
$b\mapsto\mathbf1_{K_P}\otimes b$, hence is literally $\jmath_P$.  The
middle and quotient components are $\lambda_{v,P}$ and $\alpha_{v,P}$.
Extension by zero from the edge to the global one-zero algebra completes the
bottom arrow.  All components are injective: the orbit maps are onto and
hence their function pullbacks are injective; the archimedean quotient map
sends the two minimal projections to the nonzero orthogonal projections
$e_+$ and $e_-$; and every other component is a corner embedding or a stable
isomorphism.

This proves \eqref{eq:compatible-stabilized-edge-morphism}.  Busby
functoriality for that proved morphism gives
\eqref{eq:placewise-natural-busby-equation}.  The choices are part of the
displayed coordinate model; no choice-independence or canonicality statement
is being made.
\end{proof}

Let $S_j$, $s_v=1\otimes S_{\nu(v)}\in M(J)$, and $\mu_P$ be the fixed
Cantor-pairing isometries and ideal codiagonal of
\eqref{eq:global-place-cuntz-family}--\eqref{eq:global-place-ideal-map}.
Put
\[
 T_v=1_{M(\mathfrak N_P)}\otimes S_{\nu(v)}
 \in M(\mathfrak N_P^{\mathrm s}),
 \qquad
 R_v=\overline q_P^{\mathrm s}(T_v)
 \in M(\mathfrak D_P^{\mathrm s}),
\]
where $\overline q_P^{\mathrm s}$ is the strict multiplier extension of the
quotient map.
Their restrictions to $M(\mathfrak I_P^{\mathrm s})$ satisfy
\begin{equation}
 T_v|_{M(\mathfrak I_P^{\mathrm s})}
 =\overline\jmath_P(s_v)
 \label{eq:natural-corner-compatibility}
\end{equation}
in the coordinate \eqref{eq:natural-generic-stable-coordinate}.  Define
\begin{align}
 \alpha_P((a_v)_v)
   &=\sum_{v\in V_P}R_v\alpha_{v,P}(a_v)R_v^*,
 \label{eq:stable-to-natural-quotient-map}\\
 \Lambda_P((x_v)_v)
   &=\sum_{v\in V_P}T_v\lambda_{v,P}(x_v)T_v^*.
 \label{eq:stable-to-natural-direct-sum-middle-map}
\end{align}
The sums are finite.  Orthogonality gives
\[
 R_v^*\alpha_P((a_w)_w)R_v=\alpha_{v,P}(a_v),
 \qquad
 T_v^*\Lambda_P((x_w)_w)T_v=\lambda_{v,P}(x_v),
\]
so both maps are injective.  On the direct-sum ideal,
\begin{align}
 \Lambda_P((b_v)_v)
 &=\sum_vT_v\jmath_P(b_v)T_v^*\notag\\
 &=\jmath_P\!\left(\sum_vs_vb_vs_v^*\right)
 =\jmath_P\mu_P((b_v)_v).
 \label{eq:stable-to-natural-ideal-factorization}
\end{align}

\begin{theorem}[Stable balancing extension embeds in the natural one-zero
extension]
\label{thm:stable-to-natural-extension-embedding}
For every nonempty finite prime set $P$, there is an injective morphism of
short exact sequences
\begin{equation}
\begin{CD}
0@>>>J@>>>\mathcal E_P@>>>\mathcal A_P@>>>0\\
@.@V{\jmath_P}VV@V{\Phi_P}VV@V{\alpha_P}VV@.\\
0@>>>\mathfrak I_P^{\mathrm s}
 @>>>\mathfrak N_P^{\mathrm s}
 @>{q_P^{\mathrm s}}>>\mathfrak D_P^{\mathrm s}@>>>0.
\end{CD}
\label{eq:stable-to-natural-extension-diagram}
\end{equation}
In the two Busby pullback coordinates
\eqref{eq:global-place-busby-pullback} and
\eqref{eq:natural-one-zero-busby-pullback}, its middle map is
\begin{equation}
 \Phi_P(a,m)=
 \bigl(\alpha_P(a),\overline\jmath_P(m)\bigr).
 \label{eq:stable-to-natural-middle-map}
\end{equation}
The required multiplier/corona identity is the literal equation
\begin{equation}
 \tau_P^{\mathrm{nat}}\alpha_P
 =\dot\jmath_P\tau_P.
 \label{eq:stable-to-natural-global-busby-equation}
\end{equation}
Moreover the direct-sum local middle map factors exactly:
\[
 \Lambda_P
 =\Phi_P\circ
 \left[
 ((a_v,m_v))_v\longmapsto
 \left((a_v)_v,\sum_vs_vm_vs_v^*\right)
 \right].
\]
\end{theorem}

\begin{proof}
The orthogonal corner relations and
\eqref{eq:natural-corner-compatibility} show before passing to coronas that
$\Lambda_P$ has quotient component $\alpha_P$ and ideal component
$\jmath_P\mu_P$; the latter is the literal identity
\eqref{eq:stable-to-natural-ideal-factorization}.  Hence
$(\jmath_P\mu_P,\Lambda_P,\alpha_P)$ is a morphism from the direct-sum local
extension to the natural extension.  Busby functoriality gives, for every
$a=(a_v)_v\in\mathcal A_P$,
\begin{align*}
 \tau_P^{\mathrm{nat}}\alpha_P(a)
 &=\dot{(\jmath_P\mu_P)}\,
      \widehat\tau_P^\oplus(a)\\
 &=\dot\jmath_P\dot\mu_P
      \widehat\tau_P^\oplus(a)
 =\dot\jmath_P\tau_P(a).
\end{align*}
Thus \eqref{eq:stable-to-natural-global-busby-equation} is forced by the
displayed middle, ideal, and quotient maps, not inferred from the desired
$K$-theory square.

If $(a,m)\in\mathcal E_P$, then
$\tau_P(a)=\pi_J(m)$.  The displayed Busby equation gives
\[
 \tau_P^{\mathrm{nat}}(\alpha_P(a))
 =\dot\jmath_P(\pi_J(m))
 =\pi_{\mathfrak I}(\overline\jmath_P(m)),
\]
so \eqref{eq:stable-to-natural-middle-map} lies in the natural pullback.
Coordinatewise multiplication and involution show that $\Phi_P$ is a
$*$-homomorphism.  It restricts to $\jmath_P$ on the ideal and induces
$\alpha_P$ on the quotient, proving commutativity of
\eqref{eq:stable-to-natural-extension-diagram}.

Both vertical outside maps are injective.  If $\Phi_P(e)=0$, then
$\alpha_P(q(e))=0$, hence $q(e)=0$ and $e$ lies in $J$.  Its image there is
$\jmath_P(e)=0$, so $e=0$.  Thus the middle map is injective.  Finally,
substituting the direct-sum pullback coordinates and using
\eqref{eq:stable-to-natural-ideal-factorization} proves the stated
factorization of $\Lambda_P$.
\end{proof}

The embedding has an exact inverse on its image.  In the natural pullback,
take a pair $(d,n)$ with
\[
 d\in\alpha_P(\mathcal A_P),\qquad
 n\in\overline\jmath_P(M(J)).
\]
Recover each $a_v$ by the corner compression
$R_v^*dR_v$ followed by $\alpha_{v,P}^{-1}$, and recover $m$ by applying the
multiplier extension of any evaluation map $\epsilon_{P,k}$.  Since
$\dot\jmath_P$ has the corresponding corona left inverse, the natural
pullback equation implies $\tau_P(a)=\pi_J(m)$.  Hence $(a,m)\in\mathcal
E_P$ and maps back to $(d,n)$.  Therefore
\begin{equation}
 \operatorname{im}\Phi_P
 =\left\{(d,n)\text{ in \eqref{eq:natural-one-zero-busby-pullback}}:
 d\in\operatorname{im}\alpha_P,\quad
 n\in\operatorname{im}\overline\jmath_P\right\}.
 \label{eq:stable-to-natural-image}
\end{equation}
This records the fibres and inverse data, not merely existence.

\begin{corollary}[The earlier $K$-theory square is induced by the extension
morphism]
\label{cor:stable-to-natural-K-square-lifted}
On the displayed boundary labels and ideal generator,
\[
 (\alpha_P)_*=R_P,\qquad
 (\jmath_P)_*=c_P.
\]
Consequently naturality of
\eqref{eq:stable-to-natural-extension-diagram} is exactly
\eqref{eq:stable-to-unstabilized-k-square}.
\end{corollary}

\begin{proof}
Each external corner $R_v(-)R_v^*$ induces the standard rank-one corner map
on $K$-theory.  The finite quotient maps send $[1]_p$ to $\eta_{p,P}$, and
the archimedean map sends $[p_\pm]$ to $[e_\pm]$.  This is $R_P$.
Equation \eqref{eq:constant-unit-fibre-K1-map} is $c_P$.  Naturality of
connecting maps for the proved extension morphism
\cite[Definition~21.1.1 and Example~21.1.2(a)]{blackadar1986} gives the
commuting square.
\end{proof}

There is no contradiction between injectivity of $\alpha_P$ and the fact
that $(\alpha_P)_*$ kills $[p_+]-[p_-]$.  The two projections $e_+$ and
$e_-$ remain distinct nonzero orthogonal elements of the archimedean boundary
algebra, so the algebra map from $D_2\otimes\mathcal K$ is injective.  They
are nevertheless unitarily equivalent in
$C(K_P)\rtimes C_2\cong M_2(C(H))$, so their $K_0$ classes agree.  Likewise
$\jmath_P$ is split injective, while its $K_1$ image is exactly
\[
 \mathbb Z\mathbf1_{K_P}\subset C(K_P,\mathbb Z).
\]
All nonconstant locally constant rank functions remain in the target and
outside this image.  Thus the morphism is an embedding, not an isomorphism or
a stable Morita equivalence.

For applications that require the two sign labels to remain distinct even in
$K_0$, the smallest finite-dimensional commutative symmetric split enrichment
is the one below.

\begin{proposition}[Minimal symmetric split enrichment retaining the sign
labels]
\label{prop:minimal-symmetric-sign-enrichment}
Let $\operatorname{pr}_\infty:\mathcal A_P\to\widehat A_\infty$ be the
archimedean quotient projection, and write
$q_{\mathcal E}:\mathcal E_P\to\mathcal A_P$ for the top quotient map.
Define the split-enriched natural extension
\begin{equation}
 0\to\mathfrak I_P^{\mathrm s}
 \longrightarrow
 \mathfrak N_P^{\mathrm s}\oplus\widehat A_\infty
 \longrightarrow
 \mathfrak D_P^{\mathrm s}\oplus\widehat A_\infty
 \to0,
 \label{eq:sign-enriched-natural-extension}
\end{equation}
with ideal map $i\mapsto(i,0)$ and quotient
$(x,b)\mapsto(q_P^{\mathrm s}(x),b)$.  Then
\begin{align*}
 \alpha_P^{\mathrm{sgn}}(a)
   &=(\alpha_P(a),\operatorname{pr}_\infty(a)),\\
 \Phi_P^{\mathrm{sgn}}(e)
   &=(\Phi_P(e),\operatorname{pr}_\infty(q_{\mathcal E}(e))),\\
 \jmath_P^{\mathrm{sgn}}(j)&=(\jmath_P(j),0)
\end{align*}
form an injective morphism from the stable balancing extension to
\eqref{eq:sign-enriched-natural-extension}.  Its $K_0$ map sends
\[
 [p_+]\longmapsto([e_+],[p_+]),\qquad
 [p_-]\longmapsto([e_-],[p_-]),
\]
so the sign difference survives.

Among finite-dimensional commutative split boundary carriers which retain two
nonzero orthogonal sign projections and the involution interchanging them,
$D_2=\mathbb C\oplus\mathbb C$ is minimal: every such carrier has at least
two minimal central summands, and $D_2$ has exactly two.
\end{proposition}

\begin{proof}
The second summand in \eqref{eq:sign-enriched-natural-extension} is split, so
its Busby invariant is zero.  The formulas are $*$-homomorphisms and the
previous Busby equation remains unchanged in the first component.  The second
component commutes because it is the same quotient projection on the top
middle and quotient algebras.  Injectivity follows already from $\Phi_P$.
The displayed $K_0$ images are distinct in the second direct summand.

A finite-dimensional commutative algebra is $\mathbb C^r$.  Two nonzero
orthogonal projections require disjoint nonempty supports among its $r$
minimal central projections, hence $r\geq2$.  An automorphism interchanging
the two projections permutes those supports.  The algebra
$D_2=\mathbb C\oplus\mathbb C$ with the coordinate swap realizes $r=2$,
proving minimality within the stated symmetric split class.
\end{proof}

\paragraph{Bounded coordinate replay.}
The script
\begin{center}
\footnotesize
\path{certificates/globalization_nte_stable_to_natural/}\\
\path{verify_stable_to_natural.py}
\end{center}
replays 2,361 named finite checks for one through eight finite primes.  It
passes in two independent local runtimes.  Both use Python~3.13, with distinct
SymPy versions, and verify the Cantor-corner compression identities, the constant-unit ideal factorization, the finite
matrix form of the Busby corner equation, the boundary naturality matrix, the
exact sign-difference kernel, injectivity of the split-enriched label map,
evaluation left inverses, displayed nonconstant unit-fibre classes, and the
minimal two-coordinate sign carrier.  The stable receipt ID is
\begin{center}
\footnotesize
\texttt{CERT-GNTE-STABLE-TO-}\\
\texttt{NATURAL-20260830-0001}.
\end{center}
This finite replay does not replace the essential-ideal, Green-imprimitivity, linking-algebra,
multiplier, corona, Busby-pullback, or crossed-product arguments printed
above.  It launched no Lean or Lake process.

\paragraph{Exact boundary of the result.}
The formerly missing extension-level arrow is now constructed after
stabilizing the natural one-zero extension.  No global Green reduction was
used: only the generic ideal and the individual boundary edges enter the
compatible Green modules.  Proposition
\ref{prop:unstabilized-global-action-nonproper} remains valid.  The map
$\Phi_P$ is not claimed to be onto, and
\eqref{eq:stable-to-natural-image} shows exactly what its image omits.  The
split sign enrichment retains the local character difference but does not
manufacture the nonconstant unit-fibre classes in the stable source.  No
endpoint, individual-zero, critical-line, or global zeta conclusion follows.


\subsection{The full nonconstant unit-fibre calculation: profinite orbit
coordinates, connecting maps, and the natural middle algebra}
\label{subsec:nonconstant-unit-fibre-k-theory}

The stable extension embedding of the preceding subsection sees the constant
subgroup
\(\mathbb Z\mathbf 1_{K_P}\subset C(K_P,\mathbb Z)\).  We now calculate the
remaining classes rather than referring to them collectively as
``nonconstant.''  The answer has three distinct pieces: invariant functions
on finite-place profinite quotients, coinvariants with nontrivial finite-level
transition multiplicities, and an archimedean sign-even subgroup.  All three
are needed for the K-theory of the natural one-zero algebra.

Fix a nonempty finite prime set \(P\).  For \(p\in P\), put
\begin{equation}
 Y_{p,P}=K_{P\setminus p}
 =\prod_{q\in P\setminus\{p\}}\mathbb Z_q^\times,
 \qquad
 g_{p,P}=(p)_{q\ne p}\in Y_{p,P}.
 \label{eq:unit-fibre-Y-and-g}
\end{equation}
The empty product is one point.  Let
\begin{equation}
 a_{p,P}(y)=g_{p,P}y,
 \qquad
 (\sigma_{p,P}f)(y)=f(g_{p,P}^{-1}y),
 \label{eq:unit-fibre-translation-automorphism}
\end{equation}
and let
\begin{equation}
 H_{p,P}=\overline{\langle g_{p,P}\rangle}\subset Y_{p,P},
 \qquad
 q_{p,P}:Y_{p,P}\longrightarrow Y_{p,P}/H_{p,P}
 \label{eq:unit-fibre-procyclic-closure}
\end{equation}
be the closed procyclic subgroup and quotient map.  If \(Y_{p,P}\) is not a
point, the \(\mathbb Z\)-action generated by \(a_{p,P}\) is free: an equality
\(p^ny=y\) in one factor \(\mathbb Z_q^\times\) would imply
\(p^n=1\) in \(\mathbb Q_q\), and hence as a rational number, so \(n=0\).
For \(P=\{p\}\), the action on the one-point \(Y_{p,P}\) is instead trivial;
this exceptional stabilizer is retained below and gives the ordinary circle.

\begin{lemma}[Mapping-torus model and complete finite-place K-groups]
\label{lem:unit-fibre-mapping-torus-k-groups}
There is a literal isomorphism
\begin{equation}
 C(B_{p,P})\cong M_{\sigma_{p,P}}
 =\{F\in C([0,1],C(Y_{p,P})):F(1)=\sigma_{p,P}(F(0))\}.
 \label{eq:unit-fibre-mapping-torus-isomorphism}
\end{equation}
Under this isomorphism and the Green-module coordinate for
\(\mathfrak D_{p,P}\),
\begin{align}
 K_0(\mathfrak D_{p,P})
 &\cong
 \ker(1-\sigma_{p,P,*})
 \cong C(Y_{p,P}/H_{p,P},\mathbb Z),
 \label{eq:finite-edge-K0-invariants}\\
 K_1(\mathfrak D_{p,P})
 &\cong
 \operatorname{coker}(1-\sigma_{p,P,*})
 \cong
 \frac{C(Y_{p,P},\mathbb Z)}
 {(1-\sigma_{p,P})C(Y_{p,P},\mathbb Z)}.
 \label{eq:finite-edge-K1-coinvariants}
\end{align}
No constant-function reduction occurs in either group.
\end{lemma}

\begin{proof}
Write \(\ell_p=\log p\).  For \(G\in C(B_{p,P})\), define
\[
 (\Theta G)(u)(y)=G([y,u\ell_p]),\qquad 0\leq u\leq1.
\]
Since \([y,\ell_p]=[g_{p,P}^{-1}y,0]\),
\[
 (\Theta G)(1)(y)=(\Theta G)(0)(g_{p,P}^{-1}y)
 =\sigma_{p,P}((\Theta G)(0))(y).
\]
Conversely, the endpoint equation makes a path in the displayed mapping
torus descend uniquely to the quotient
\((Y_{p,P}\times\mathbb R)/\langle\beta_p\rangle\).  These two maps preserve
multiplication, involution, and the supremum norm, proving
\eqref{eq:unit-fibre-mapping-torus-isomorphism}.

We also record the coefficient groups rather than hiding them behind a
topological K-theory identification.  For every compact profinite space
\(Y\), a projection over \(Y\) is, after refining a finite clopen partition,
unitarily equivalent to a constant matrix projection on each cell.  Stable
equivalence is therefore determined exactly by its locally constant rank
function, so
\[
 K_0(C(Y))=C(Y,\mathbb Z).
\]
Moreover every continuous unitary matrix over \(Y\) is uniformly approximated
within distance less than two by one that is constant on a finite clopen
partition.  The two unitaries are homotopic, and each constant value is
homotopic to the identity inside a unitary group.  Thus
\(K_1(C(Y))=0\).

Blackadar's mapping-torus exact sequence and its explicit connecting maps are
\(1-\sigma_*\) under the Bott coordinate
\cite[Theorem~10.2.1, Definition~10.3.1, and Proposition~10.4.1,
pp.~83--87]{blackadar1986}.  Substitution of the two coefficient groups just
proved gives the kernel and cokernel in
\eqref{eq:finite-edge-K0-invariants}--
\eqref{eq:finite-edge-K1-coinvariants}.  A locally constant function is fixed
by \(\sigma_{p,P}\) exactly when it is constant on every orbit of
\(g_{p,P}\), hence on every coset of the closure \(H_{p,P}\).  Pullback along
\(q_{p,P}\) therefore identifies the kernel with
\(C(Y_{p,P}/H_{p,P},\mathbb Z)\).  Finally, the restricted Green module from
Theorem~\ref{thm:multiplace-finite-edge-coordinates} transports these groups
to \(\mathfrak D_{p,P}\) without changing the extension boundary coordinate.
\end{proof}

The coinvariant group in \eqref{eq:finite-edge-K1-coinvariants} generally
does not reduce to a free group on a final orbit set.  Its exact finite-level
coordinates are as follows.  Let \(\mathcal M_{p,P}\) be the directed set of
tuples \(\mathbf m=(m_q)_{q\ne p}\), with \(m_q\geq1\), ordered
coordinatewise.  For an empty tuple set all objects below equal to the
one-point object and put \(L_{\varnothing}=1\).  Otherwise define
\begin{equation}
 G_{\mathbf m}=\prod_{q\ne p}
 (\mathbb Z/q^{m_q}\mathbb Z)^\times,
 \quad
 g_{\mathbf m}=(p\bmod q^{m_q})_{q\ne p},
 \quad
 L_{\mathbf m}=\operatorname{ord}_{G_{\mathbf m}}(g_{\mathbf m}),
 \quad
 \Omega_{\mathbf m}=G_{\mathbf m}/\langle g_{\mathbf m}\rangle.
 \label{eq:finite-orbit-coordinate-data}
\end{equation}
Every orbit has exactly \(L_{\mathbf m}\) points.  Thus the stabilizer of a
point for the finite-level \(\mathbb Z\)-action is precisely
\(L_{\mathbf m}\mathbb Z\); these finite-level stabilizers disappear in the
nontrivial inverse limit but are not deleted from the transition maps.

For an orbit \(O\in\Omega_{\mathbf m}\), let
\[
 u_O=\mathbf1_O
 \quad\hbox{in the invariant group},
 \qquad
 v_O=[\delta_x]\quad(x\in O)
 \quad\hbox{in the coinvariant group}.
\]
The second class is independent of \(x\in O\).  If
\(\mathbf m'\geq\mathbf m\), let
\(\pi_{\mathbf m',\mathbf m}:G_{\mathbf m'}\to G_{\mathbf m}\) be reduction.
Then \(L_{\mathbf m}\mid L_{\mathbf m'}\), and the two transition maps are
\begin{align}
 \iota^0_{\mathbf m',\mathbf m}(u_O)
 &=\sum_{\substack{O'\in\Omega_{\mathbf m'}\\
                    O'\subset\pi_{\mathbf m',\mathbf m}^{-1}(O)}}u_{O'},
 \label{eq:finite-invariant-transition}\\
 \iota^1_{\mathbf m',\mathbf m}(v_O)
 &=\frac{L_{\mathbf m'}}{L_{\mathbf m}}
   \sum_{\substack{O'\in\Omega_{\mathbf m'}\\
                    O'\subset\pi_{\mathbf m',\mathbf m}^{-1}(O)}}v_{O'}.
 \label{eq:finite-coinvariant-transition}
\end{align}

\begin{proposition}[Finite-level coordinate form, including multiplicities]
\label{prop:unit-fibre-finite-level-limits}
The two groups in
\eqref{eq:finite-edge-K0-invariants}--
\eqref{eq:finite-edge-K1-coinvariants} are respectively
\begin{align}
 C(Y_{p,P}/H_{p,P},\mathbb Z)
 &\cong
 \varinjlim_{\mathbf m}
 (\mathbb Z[\Omega_{\mathbf m}],\iota^0),
 \label{eq:finite-invariant-direct-limit}\\
 C(Y_{p,P},\mathbb Z)_{\sigma_{p,P}}
 &\cong
 \varinjlim_{\mathbf m}
 (\mathbb Z[\Omega_{\mathbf m}],\iota^1).
 \label{eq:finite-coinvariant-direct-limit}
\end{align}
Both displayed transition maps are injective.  In particular the
coinvariant group is torsion-free, but no free decomposition or canonical
splitting is asserted.
\end{proposition}

\begin{proof}
Every locally constant integer-valued function on \(Y_{p,P}\) factors through
some \(G_{\mathbf m}\), and pullback along the reduction maps gives
\[
 C(Y_{p,P},\mathbb Z)
 =\varinjlim_{\mathbf m}\mathbb Z^{G_{\mathbf m}}.
\]
At level \(\mathbf m\), the invariant functions have basis
\((u_O)_{O\in\Omega_{\mathbf m}}\), and pullback of \(\mathbf1_O\) is the
sum of the characteristic functions of all fine orbits over \(O\).  This is
\eqref{eq:finite-invariant-transition}.

The coinvariants of the permutation module of one orbit are \(\mathbb Z\),
with every point mass representing the same generator \(v_O\).  Fix
\(x\in O\).  Pullback of \(\delta_x\) is the sum of the point masses over
\(\pi^{-1}(x)\).  Every fine orbit \(O'\) above \(O\) maps onto \(O\) with
degree \(L_{\mathbf m'}/L_{\mathbf m}\), so exactly that many of its points
lie above \(x\).  Hence its contribution is
\((L_{\mathbf m'}/L_{\mathbf m})v_{O'}\), proving
\eqref{eq:finite-coinvariant-transition}.  Equivalently, the number of fine
orbits above \(O\) is
\[
 |\ker\pi_{\mathbf m',\mathbf m}|
 \frac{L_{\mathbf m}}{L_{\mathbf m'}},
\]
which verifies that all points in the pullback have been counted exactly
once.

Filtered colimits commute with the cokernel of
\(1-\sigma_{p,P}\), giving
\eqref{eq:finite-coinvariant-direct-limit}.  For invariants, a function that
factors through \(G_{\mathbf m}\) is invariant in the inverse limit exactly
when its finite-level function is constant on the
\(\langle g_{\mathbf m}\rangle\)-orbits, which gives
\eqref{eq:finite-invariant-direct-limit}.  Distinct coarse orbits have
disjoint inverse images.  Both transition formulas therefore send a nonzero
integer combination to a nonzero combination on disjoint supports.  They are
injective.  An injective filtered union of torsion-free groups is
torsion-free, proving the final statement.
\end{proof}

We next compute the connecting map on every invariant rank function.  Let
\begin{equation}
 \operatorname{pr}_{p,P}:K_P\longrightarrow Y_{p,P}
 \label{eq:finite-unit-drop-map}
\end{equation}
drop the \(p\)-coordinate, and define
\begin{equation}
 \iota_{p,P}:C(Y_{p,P}/H_{p,P},\mathbb Z)
 \longrightarrow C(K_P,\mathbb Z),
 \qquad
 \iota_{p,P}(h)=h\circ q_{p,P}\circ\operatorname{pr}_{p,P}.
 \label{eq:finite-invariant-ideal-pullback}
\end{equation}

\begin{theorem}[Full finite-edge connecting map and middle K-groups]
\label{thm:full-finite-edge-boundary-map}
In the coordinates above, the connecting map of the \(p\)-edge is
\begin{equation}
 \partial_{p,P}(h)=-\iota_{p,P}(h)
 \qquad
 (h\in C(Y_{p,P}/H_{p,P},\mathbb Z)).
 \label{eq:full-finite-edge-boundary-map}
\end{equation}
It is injective.  Consequently
\begin{equation}
 K_0(A_{\varnothing p}^{(P)})=0
 \label{eq:finite-edge-middle-K0-zero}
\end{equation}
and there is an exact sequence
\begin{equation}
 0\longrightarrow
 \frac{C(K_P,\mathbb Z)}
      {\iota_{p,P}C(Y_{p,P}/H_{p,P},\mathbb Z)}
 \longrightarrow K_1(A_{\varnothing p}^{(P)})
 \longrightarrow C(Y_{p,P},\mathbb Z)_{\sigma_{p,P}}
 \longrightarrow0.
 \label{eq:finite-edge-middle-K1-exact}
\end{equation}
No splitting of \eqref{eq:finite-edge-middle-K1-exact} is selected.
\end{theorem}

\begin{proof}
It is enough by additivity to take
\(h=\mathbf1_E\) for a clopen subset
\(E\subset Y_{p,P}/H_{p,P}\).  Put
\(\widetilde E=q_{p,P}^{-1}(E)\).  It is clopen and invariant under
\(g_{p,P}\).  The corresponding clopen part of the mapping-torus boundary is
\[
 B_E=(\widetilde E\times\mathbb R)/\langle\beta_p\rangle
 \subset B_{p,P}.
\]
In the orbit coordinates of
Theorem~\ref{thm:multiplace-finite-edge-coordinates}, the scalar function
\begin{equation}
 x_E(a,t)=
 \frac{\mathbf1_{\widetilde E}(\operatorname{pr}_{p,P}(a))}{1+e^t},
 \qquad
 x_E([b,s])=\mathbf1_{\widetilde E}(b)
 \label{eq:finite-edge-explicit-lift}
\end{equation}
is continuous.  Invariance of \(\widetilde E\) makes the boundary formula
independent of the representative \((b,s)\), while the exact cross-convergence
criterion makes the generic value tend to it as \(t\to-\infty\).  The only
escaping end is \(t\to+\infty\), where the function vanishes.  Thus \(x_E\)
is a self-adjoint lift of the boundary projection \(\mathbf1_{B_E}\).

Its exponential is the identity off
\(\operatorname{pr}_{p,P}^{-1}(\widetilde E)\).  On every retained generic
unit fibre above \(\widetilde E\), its argument decreases from \(2\pi\) to
zero as \(t\) increases, so it has winding \(-1\) in the previously declared
orientation.  Hence
\[
 \partial_{p,P}[\mathbf1_{B_E}]
 =-\mathbf1_{\operatorname{pr}_{p,P}^{-1}(\widetilde E)}.
\]
Every locally constant integer function is an integer combination of such
characteristic functions, proving
\eqref{eq:full-finite-edge-boundary-map}.  The map
\(q_{p,P}\circ\operatorname{pr}_{p,P}\) is onto, so its pullback is injective.

Finally, \(\mathfrak I_P\) is Morita equivalent to
\(C_0(K_P\times\mathbb R)\), whose K-groups in the retained suspension
coordinate are
\[
 K_0(\mathfrak I_P)=0,
 \qquad K_1(\mathfrak I_P)=C(K_P,\mathbb Z).
\]
Insert these, Lemma~\ref{lem:unit-fibre-mapping-torus-k-groups}, and the
injective connecting map into the edge six-term sequence.  Exactness gives
\eqref{eq:finite-edge-middle-K0-zero} and
\eqref{eq:finite-edge-middle-K1-exact}.
\end{proof}

The archimedean summand contributes all sign-even functions, not only its
constant unit.  Retain the clopen half \(H_P^+\subset K_P\) constructed in
the proof of Theorem~\ref{thm:multiplace-arch-boundary}, so
\begin{equation}
 K_P=H_P^+\sqcup(-H_P^+).
 \label{eq:arch-clopen-half}
\end{equation}
Define the continuous representative map
\begin{equation}
 r_P:K_P\longrightarrow H_P^+,
 \qquad
 r_P(c)=
 \begin{cases}c,&c\in H_P^+,\\-c,&c\in-H_P^+.
 \end{cases}
 \label{eq:arch-representative-map}
\end{equation}

\begin{theorem}[Full archimedean boundary map]
\label{thm:full-archimedean-unit-fibre-boundary}
The explicit matrix coordinate gives
\begin{equation}
 K_0(\mathfrak D_{\infty,P})\cong C(H_P^+,\mathbb Z),
 \qquad
 K_1(\mathfrak D_{\infty,P})=0.
 \label{eq:arch-boundary-full-K-groups}
\end{equation}
Under the generic coordinate
\(K_1(\mathfrak I_P)=C(K_P,\mathbb Z)\), the complete archimedean
connecting map is
\begin{equation}
 \partial_{\infty,P}(h)=-r_P^*h.
 \label{eq:full-archimedean-boundary-map}
\end{equation}
Its image is exactly the sign-even subgroup
\begin{equation}
 C(K_P,\mathbb Z)^{\mathrm{ev}}
 =\{f:f(c)=f(-c)\text{ for all }c\in K_P\}.
 \label{eq:sign-even-unit-functions}
\end{equation}
\end{theorem}

\begin{proof}
The clopen decomposition \eqref{eq:arch-clopen-half} trivializes the free
double cover, and the already displayed orbit matrices give
\[
 C(K_P)\rtimes C_2\cong M_2(C(H_P^+)).
\]
The profinite K-theory calculation in the proof of
Lemma~\ref{lem:unit-fibre-mapping-torus-k-groups} yields
\eqref{eq:arch-boundary-full-K-groups}.

For the sign and orientation, write every \(c\in K_P\) uniquely as
\(c=\varepsilon h\), with \(h\in H_P^+\) and
\(\varepsilon\in\{1,-1\}\).  The orbit map for the diagonal action on
\(K_P\times\mathbb R\) is explicitly
\begin{equation}
 [(\varepsilon h,z)]\longmapsto(h,w=\varepsilon z)
 \quad\hbox{in }H_P^+\times\mathbb R.
 \label{eq:arch-orbit-coordinate-w}
\end{equation}
On the generic part \(w\ne0\), its relation with the retained ideal coordinate
\((a,t)\in K_P\times\mathbb R\) is
\begin{equation}
 a=\operatorname{sgn}(w)h,
 \qquad t=\log|w|.
 \label{eq:arch-w-to-generic-coordinate}
\end{equation}

For a clopen \(E\subset H_P^+\), the function
\[
 x_E(h,w)=\frac{\mathbf1_E(h)}{1+|w|}
\]
on the commutative orbit extension is a self-adjoint lift of
\(\mathbf1_E\) at \(w=0\).  On both components \(w>0\) and \(w<0\), the
argument of its exponential decreases from \(2\pi\) to zero as
\(t=\log|w|\) increases.  Formula
\eqref{eq:arch-w-to-generic-coordinate} therefore gives pointwise winding
\(-1\) on
\(E\sqcup(-E)=r_P^{-1}(E)\) and zero elsewhere.  The extension-level Green
module transports this calculation to the crossed-product extension, so
\[
 \partial_{\infty,P}[\mathbf1_E]
 =-\mathbf1_{r_P^{-1}(E)}.
\]
Additivity proves \eqref{eq:full-archimedean-boundary-map}.  Pullback by
\(r_P\) is injective and its image consists exactly of the functions constant
on each pair \(\{c,-c\}\), proving
\eqref{eq:sign-even-unit-functions}.
\end{proof}

We can now assemble all places.  Abbreviate
\begin{equation}
 \mathcal F_{p,P}=C(Y_{p,P}/H_{p,P},\mathbb Z),
 \qquad
 \mathcal C_{p,P}=C(Y_{p,P},\mathbb Z)_{\sigma_{p,P}}.
 \label{eq:finite-invariant-coinvariant-abbreviations}
\end{equation}

\begin{theorem}[Full K-theory and connecting map of the natural one-zero
extension]
\label{thm:full-natural-one-zero-k-theory}
The quotient groups are
\begin{align}
 K_0(\mathfrak D_P)
 &\cong
 \left(\bigoplus_{p\in P}\mathcal F_{p,P}\right)
 \oplus C(H_P^+,\mathbb Z),
 \label{eq:global-natural-boundary-K0}\\
 K_1(\mathfrak D_P)
 &\cong\bigoplus_{p\in P}\mathcal C_{p,P}.
 \label{eq:global-natural-boundary-K1}
\end{align}
In these exact coordinates the connecting map is
\begin{equation}
 \partial_P^{(1)}((h_p)_{p\in P},h_\infty)
 =-\sum_{p\in P}\iota_{p,P}(h_p)-r_P^*h_\infty
 \quad\text{in }C(K_P,\mathbb Z).
 \label{eq:global-full-unit-fibre-boundary-map}
\end{equation}

For a finite tuple \(h=(h_p)_{p\in P}\), put
\begin{equation}
 S_h(c)=\sum_{p\in P}\iota_{p,P}(h_p)(c).
 \label{eq:finite-place-sum-function}
\end{equation}
Then projection to the finite-place coordinates gives an isomorphism
\begin{equation}
 K_0(\mathfrak N_P)
 \xrightarrow{\ \cong\ }
 \left\{h\in\bigoplus_{p\in P}\mathcal F_{p,P}:
 S_h(c)=S_h(-c)\text{ for every }c\in K_P\right\},
 \label{eq:natural-middle-K0-exact-kernel}
\end{equation}
whose inverse is
\begin{equation}
 h\longmapsto((h_p)_p,-S_h|_{H_P^+}).
 \label{eq:natural-middle-K0-kernel-inverse}
\end{equation}

Define
\begin{equation}
 d_{p,P}:\mathcal F_{p,P}\longrightarrow C(H_P^+,\mathbb Z),
 \qquad
 d_{p,P}(h)(c)=\iota_{p,P}(h)(c)-\iota_{p,P}(h)(-c),
 \label{eq:finite-place-odd-difference-map}
\end{equation}
and
\begin{equation}
 \mathcal Q_P=
 \frac{C(H_P^+,\mathbb Z)}
 {\sum_{p\in P}d_{p,P}(\mathcal F_{p,P})}.
 \label{eq:global-unit-fibre-cokernel-coordinate}
\end{equation}
There is a short exact sequence
\begin{equation}
 0\longrightarrow\mathcal Q_P
 \longrightarrow K_1(\mathfrak N_P)
 \longrightarrow\bigoplus_{p\in P}\mathcal C_{p,P}
 \longrightarrow0.
 \label{eq:natural-middle-K1-short-exact}
\end{equation}
No canonical splitting of \eqref{eq:natural-middle-K1-short-exact} is
asserted.
\end{theorem}

\begin{proof}
The quotient \(\mathfrak D_P\) is the finite direct sum of its boundary
strata.  Lemma~\ref{lem:unit-fibre-mapping-torus-k-groups} and
Theorem~\ref{thm:full-archimedean-unit-fibre-boundary} therefore give
\eqref{eq:global-natural-boundary-K0} and
\eqref{eq:global-natural-boundary-K1}.  Extension by zero from every edge to
the global one-zero algebra is a morphism of extensions.  Naturality and
additivity of the connecting maps, now applied to the complete edge groups
rather than only their units, combine
\eqref{eq:full-finite-edge-boundary-map} and
\eqref{eq:full-archimedean-boundary-map} into
\eqref{eq:global-full-unit-fibre-boundary-map}.

Since \(K_0(\mathfrak I_P)=0\) and
\(K_1(\mathfrak I_P)=C(K_P,\mathbb Z)\), the six-term sequence gives
\[
 K_0(\mathfrak N_P)=\ker\partial_P^{(1)}
\]
and
\begin{equation}
 0\to\operatorname{coker}\partial_P^{(1)}
 \to K_1(\mathfrak N_P)
 \to K_1(\mathfrak D_P)\to0.
 \label{eq:natural-six-term-middle-fragment}
\end{equation}
Equation \eqref{eq:global-full-unit-fibre-boundary-map} vanishes exactly when
\(S_h=-r_P^*h_\infty\).  The right side is sign-even.  Thus \(S_h\) must be
sign-even; when it is, the unique solution is
\(h_\infty=-S_h|_{H_P^+}\).  This proves
\eqref{eq:natural-middle-K0-exact-kernel} and its explicit inverse.

It remains to identify the cokernel without discarding the archimedean image.
Define
\begin{equation}
 \operatorname{Odd}_{H_P^+}:C(K_P,\mathbb Z)
 \longrightarrow C(H_P^+,\mathbb Z),
 \qquad
 \operatorname{Odd}_{H_P^+}(f)(c)=f(c)-f(-c).
 \label{eq:odd-difference-quotient-map}
\end{equation}
This map is onto: prescribe a function on \(H_P^+\) and extend it by zero on
\(-H_P^+\).  Its kernel is exactly the sign-even subgroup
\(r_P^*C(H_P^+,\mathbb Z)\).  Hence it induces the exact isomorphism
\[
 \frac{C(K_P,\mathbb Z)}{r_P^*C(H_P^+,\mathbb Z)}
 \xrightarrow{\ \cong\ }C(H_P^+,\mathbb Z).
\]
Under this map, the finite-place summand \(\iota_{p,P}(h)\) becomes exactly
\(d_{p,P}(h)\).  Therefore
\(\operatorname{coker}\partial_P^{(1)}\cong\mathcal Q_P\).  Substitution into
\eqref{eq:natural-six-term-middle-fragment} proves
\eqref{eq:natural-middle-K1-short-exact}.
\end{proof}

\paragraph{The one-prime coordinate check.}
If \(P=\{p\}\), then \(Y_{p,P}\), \(H_{p,P}\), and every finite quotient
are one point.  Thus
\(\mathcal F_{p,P}=\mathcal C_{p,P}=\mathbb Z\).  The finite edge has
\[
 K_0(A_{\varnothing p}^{(P)})=0,
 \qquad
 0\to C(K_P,\mathbb Z)/\mathbb Z\mathbf1_{K_P}
 \to K_1(A_{\varnothing p}^{(P)})\to\mathbb Z\to0.
\]
For the global natural algebra, the K-zero kernel consists exactly of
\((n,-n\mathbf1_{H_P^+})\), so \(K_0(\mathfrak N_P)\cong\mathbb Z\).  The
finite odd-difference map is zero, and therefore
\[
 0\to C(H_P^+,\mathbb Z)\to K_1(\mathfrak N_P)\to\mathbb Z\to0.
\]
This last sequence splits abstractly because its quotient is \(\mathbb Z\),
but the calculation selects no splitting.  The term
\(C(H_P^+,\mathbb Z)\) is the exact residual sign-asymmetric unit-fibre
coordinate after the archimedean boundary has filled every sign-even function.

\paragraph{Exact comparison with the stable balancing extension.}
For constant \(h_p=k_p\) and constant \(h_\infty=n\),
\eqref{eq:global-full-unit-fibre-boundary-map} reduces to
\[
 -\left(\sum_{p\in P}k_p+n\right)\mathbf1_{K_P}.
\]
The two stable archimedean character projections both map to the same
constant \(n=1\) class, exactly as proved in
Corollary~\ref{cor:stable-to-natural-K-square-lifted}.  Thus the earlier
codiagonal square is the constant-coordinate restriction of the full map,
not a substitute for it.  The new calculation also shows precisely what lies
outside that image: finite invariant functions, finite coinvariants with the
multiplicities \eqref{eq:finite-coinvariant-transition}, and the odd-difference
quotient \(\mathcal Q_P\).  It proves no surjectivity, Morita equivalence,
choice-independence across prime sets, endpoint statement, or global zeta
consequence.

\paragraph{Bounded finite-coordinate replay.}
The executable certificate and its receipt are stored under
\begin{center}
\footnotesize
\path{certificates/globalization_nte_nonconstant_unit_k/}.
\end{center}
The receipt has stable ID
\texttt{CERT-GNTE-NONCONSTANT-UNIT-K-20260830-0001}.
They verify the finite quotient orbit lengths, finite-level stabilizers, invariant
and coinvariant bases, both transition formulas including their
\(L_{\mathbf m'}/L_{\mathbf m}\) multiplicities, injectivity of the
transition matrices, the finite and archimedean pullback matrices, and the
global kernel/cokernel coordinate identities on bounded examples.  The
certificate does not replace the mapping-torus theorem, Green imprimitivity,
six-term exactness, or the proofs above, and it launches no Lean or Lake
process.


\subsection{Finite-set coherence: exact tail coordinates, transition morphisms,
and the multiplier obstruction}
\label{subsec:finite-set-coherence}

The preceding stable-to-natural theorem was proved for one finite prime set
$P$ at a time.  Its stabilization coordinate $\chi_P$ was explicitly declared
to be a choice.  Equality of the formulas for two different sets does not
prove that the chosen isomorphisms commute when a prime is adjoined.  We now
construct one simultaneous coordinate system and prove every transition
square in that system.  The construction also exposes a genuine obstruction:
when $P\subsetneq P'$, the natural middle-algebra map occupies the sector in
which every new finite coordinate is nonzero.  It is therefore degenerate as
a map into the full $P'$ algebra and has no unital strict multiplier extension
there.  The generic-ideal map is nevertheless nondegenerate, so the multiplier
and corona maps required for the Busby square do exist and are explicit.

Let $\mathbb P$ denote the set of rational primes.  For a finite prime set
$P$ define the positive complementary prime group
\begin{equation}
 \mathcal L_P=
 \left\{\prod_{r\in\mathbb P\setminus P}r^{n_r}:
 n_r\in\mathbb Z,\ n_r=0\text{ for all but finitely many }r\right\}.
 \label{eq:coherence-complementary-prime-group}
\end{equation}
If $P\subset P'$, put
\begin{equation}
 R=P'\setminus P,\qquad
 \Gamma_R^+=\left\{\prod_{r\in R}r^{n_r}:n_r\in\mathbb Z\right\},
 \qquad K_R=\prod_{r\in R}\mathbb Z_r^\times.
 \label{eq:coherence-new-prime-data}
\end{equation}
Unique prime factorization gives literal direct products
\begin{equation}
 \Gamma_{P'}=\Gamma_P\times\Gamma_R^+,
 \qquad
 \mathcal L_P=\Gamma_R^+\times\mathcal L_{P'},
 \qquad
 \Gamma_P\times\mathcal L_P=\mathbb Q^\times.
 \label{eq:coherence-group-factorizations}
\end{equation}
The sign belongs to $\Gamma_P$; the other two groups in the first two
factorizations contain only positive rationals.  Hence no sign coordinate is
hidden in these equations.

\subsubsection{The old-boundary sector and its complete coordinates}

Inside $W_{P'}^{(1)}$ define the invariant open subset
\begin{equation}
 U_{P,P'}=\{x\in W_{P'}^{(1)}:x_r\ne0\text{ for every }r\in R\}.
 \label{eq:coherence-old-boundary-sector}
\end{equation}
It contains the generic stratum and exactly the one-zero strata labelled by
$P\sqcup\{\infty\}$; the new-prime boundary strata are not identified with
those old strata.

\begin{theorem}[Coordinate homeomorphism retaining every new unit]
\label{thm:coherence-sector-homeomorphism}
For $x=((x_q)_{q\in P'},y)\in U_{P,P'}$, write
\[
 m_r=v_r(x_r),\qquad u_r=r^{-m_r}x_r\in\mathbb Z_r^\times
 \quad(r\in R),
 \qquad h_R(x)=\prod_{r\in R}r^{m_r}\in\Gamma_R^+.
\]
Define
\begin{align}
 z_q&=h_R(x)^{-1}x_q &&(q\in P),
 &z_\infty&=h_R(x)^{-1}y,
 \label{eq:coherence-old-coordinate}\\
 b_r&=h_R(x)^{-1}x_r
 =\left(\prod_{s\in R\setminus\{r\}}s^{-m_s}\right)u_r
 \in\mathbb Z_r^\times &&(r\in R).
 \label{eq:coherence-new-unit-coordinate}
\end{align}
Then
\begin{equation}
 \Theta_{P,P'}:U_{P,P'}\xrightarrow{\ \cong\ }
 \Gamma_R^+\times\bigl(W_P^{(1)}\times K_R\bigr),
 \qquad x\longmapsto\bigl(h_R(x),z,(b_r)_{r\in R}\bigr)
 \label{eq:coherence-sector-homeomorphism}
\end{equation}
is a homeomorphism.  Its inverse is the coordinatewise formula
\begin{equation}
 x_q=h z_q\ (q\in P),\qquad
 x_r=h b_r\ (r\in R),\qquad y=h z_\infty.
 \label{eq:coherence-sector-homeomorphism-inverse}
\end{equation}
For $\gamma\in\Gamma_P$ and $\rho\in\Gamma_R^+$, the action is
\begin{equation}
 \Theta_{P,P'}((\gamma\rho)x)
 =\bigl(\rho h_R(x),\gamma z,(\gamma b_r)_{r\in R}\bigr).
 \label{eq:coherence-sector-action}
\end{equation}
Thus $\Gamma_R^+$ translates only the first discrete factor, whereas
$\Gamma_P$ acts diagonally on $W_P^{(1)}\times K_R$.
\end{theorem}

\begin{proof}
Since $v_r(h_R(x))=m_r$, equation
\eqref{eq:coherence-new-unit-coordinate} has $r$-adic valuation zero.
Every factor $s^{-m_s}$ with $s\ne r$ is an $r$-adic unit, so the expanded
formula retains all cross-prime factors and proves $b_r\in\mathbb Z_r^\times$.
Multiplication by the nonzero rational $h_R(x)^{-1}$ does not change which of
the old coordinates is zero.  Hence $z\in W_P^{(1)}$.

Given $(h,z,b)$ on the right, write $h=\prod_{r\in R}r^{m_r}$ and use
\eqref{eq:coherence-sector-homeomorphism-inverse}.  Its new $r$ coordinate
has valuation $m_r$, and applying
\eqref{eq:coherence-new-unit-coordinate} returns $b_r$.  Applying
\eqref{eq:coherence-old-coordinate} returns $z$.  These calculations prove
both inverse identities.

On $\mathbb Q_r^\times$, the valuation is locally constant and the unit map
$x_r\mapsto r^{-v_r(x_r)}x_r$ is continuous.  The group
$\Gamma_R^+$ is discrete.  Therefore $\Theta_{P,P'}$ is continuous on each
clopen valuation sector, and the inverse formula is continuous.  It is a
homeomorphism.

Finally, $\gamma$ has zero $r$-adic valuation for every $r\in R$, whereas
$v_r(\rho x_r)=v_r(\rho)+m_r$.  Thus
$h_R((\gamma\rho)x)=\rho h_R(x)$.  Substitution cancels $\rho$ in every
old and new unit coordinate and leaves multiplication by $\gamma$, proving
\eqref{eq:coherence-sector-action}.
\end{proof}

The cross-prime factors in
\eqref{eq:coherence-new-unit-coordinate} are essential.  Replacing $b_r$ by
$u_r$ would make $\Gamma_R^+$ alter the alleged base coordinate and would
destroy both the product action and the triple-inclusion calculation below.

Put
\[
 \mathfrak B_{P,R}
 =C_0(W_P^{(1)}\times K_R)\rtimes\Gamma_P.
\]
Projection to $W_P^{(1)}$ is proper and equivariant, because $K_R$ is
compact.  It therefore gives a nondegenerate injective map
\begin{equation}
 \operatorname{pr}_{P,R}^{\rtimes}:\mathfrak N_P\longrightarrow
 \mathfrak B_{P,R},\qquad
 fu_\gamma\longmapsto(f\circ\operatorname{pr}_{P,R})u_\gamma.
 \label{eq:coherence-projection-crossed-product-map}
\end{equation}
Injectivity follows directly from the faithful conditional expectations:
if the image of $a$ is zero, then the image of $E(a^*a)$ is zero; coefficient
pullback is injective, so $E(a^*a)=0$ and $a=0$.  We use reduced crossed
products in this argument and then full equals reduced because all groups are
amenable.

Let $e_{\rho,\sigma}$ be the matrix units on
$\ell^2(\Gamma_R^+)$.  The homeomorphism above gives a crossed-product
isomorphism
\begin{equation}
 \Xi_{P,P'}:
 \mathcal K(\ell^2\Gamma_R^+)\otimes\mathfrak B_{P,R}
 \xrightarrow{\ \cong\ }
 C_0(U_{P,P'})\rtimes\Gamma_{P'}.
 \label{eq:coherence-sector-crossed-product-isomorphism}
\end{equation}
On the dense algebraic span, with the action convention of
\eqref{eq:coherence-sector-action}, it is exactly
\begin{equation}
 e_{\rho,\sigma}\otimes(fu_\gamma)
 \longmapsto
 (\mathbf1_{\{\rho\}}\otimes f)
 u_{(\rho\sigma^{-1},\gamma)}.
 \label{eq:coherence-sector-matrix-units}
\end{equation}
Multiplying two displayed generators gives
$\delta_{\sigma,\rho'}e_{\rho,\sigma'}$ on both sides, and involution
interchanges the two matrix indices.  The regular covariant representation
identifies $C_0(\Gamma_R^+)\rtimes\Gamma_R^+$ with
$\mathcal K(\ell^2\Gamma_R^+)$, so the dense formula extends to the asserted
isomorphism.

\subsubsection{Tail stabilization and the natural transition extension}

Fix $H_0=\ell^2(\mathbb N_0)$ and put
\begin{align}
 \widetilde{\mathfrak N}_P
 &=\mathfrak N_P\otimes\mathcal K(\ell^2\mathcal L_P)
                  \otimes\mathcal K(H_0),\notag\\
 \widetilde{\mathfrak I}_P
 &=\mathfrak I_P\otimes\mathcal K(\ell^2\mathcal L_P)
                  \otimes\mathcal K(H_0),\notag\\
 \widetilde{\mathfrak D}_P
 &=\mathfrak D_P\otimes\mathcal K(\ell^2\mathcal L_P)
                  \otimes\mathcal K(H_0).
 \label{eq:coherence-tail-stabilized-natural-algebras}
\end{align}
These form the exact tail-stabilized natural extension.  For
$P\subset P'$, the unique factorization
$\ell=\rho\ell'$ in
$\mathcal L_P=\Gamma_R^+\times\mathcal L_{P'}$ gives the unitary
\begin{equation}
 V_{P,P'}:\ell^2(\mathcal L_P)\longrightarrow
 \ell^2(\Gamma_R^+)\otimes\ell^2(\mathcal L_{P'}),
 \qquad \delta_{\rho\ell'}\longmapsto\delta_\rho\otimes\delta_{\ell'}.
 \label{eq:coherence-tail-unitary}
\end{equation}

Extension by zero from the invariant open set $U_{P,P'}$, together with
\eqref{eq:coherence-projection-crossed-product-map}--
\eqref{eq:coherence-tail-unitary}, defines an injective homomorphism
\begin{equation}
 \Psi_{P,P'}:\widetilde{\mathfrak N}_P
 \longrightarrow\widetilde{\mathfrak N}_{P'}.
 \label{eq:coherence-natural-middle-transition}
\end{equation}
Written as a composition with every carrier visible, it is
\begin{align}
 \mathfrak N_P\otimes\mathcal K(\ell^2\mathcal L_P)\otimes\mathcal K(H_0)
 &\xrightarrow{\operatorname{Ad}V_{P,P'}}
 \mathcal K(\ell^2\Gamma_R^+)\otimes\mathfrak N_P
 \otimes\mathcal K(\ell^2\mathcal L_{P'})\otimes\mathcal K(H_0)
 \notag\\
 &\xrightarrow{1\otimes\operatorname{pr}_{P,R}^{\rtimes}\otimes1}
 \mathcal K(\ell^2\Gamma_R^+)\otimes\mathfrak B_{P,R}
 \otimes\mathcal K(\ell^2\mathcal L_{P'})\otimes\mathcal K(H_0)
 \notag\\
 &\xrightarrow{\Xi_{P,P'}\otimes1}
 (C_0(U_{P,P'})\rtimes\Gamma_{P'})
 \otimes\mathcal K(\ell^2\mathcal L_{P'})\otimes\mathcal K(H_0)
 \notag\\
 &\xrightarrow{\operatorname{ext}_{U}\otimes1}
 \widetilde{\mathfrak N}_{P'}.
 \label{eq:coherence-natural-middle-transition-expanded}
\end{align}
It restricts to an ideal map $\Psi^I_{P,P'}$ and induces a quotient map
$\Psi^D_{P,P'}$, giving a morphism
\begin{equation}
\begin{CD}
0@>>>\widetilde{\mathfrak I}_P@>>>\widetilde{\mathfrak N}_P
 @>>>\widetilde{\mathfrak D}_P@>>>0\\
@.@V{\Psi^I_{P,P'}}VV@V{\Psi_{P,P'}}VV@V{\Psi^D_{P,P'}}VV@.\\
0@>>>\widetilde{\mathfrak I}_{P'}@>>>\widetilde{\mathfrak N}_{P'}
 @>>>\widetilde{\mathfrak D}_{P'}@>>>0.
\end{CD}
\label{eq:coherence-natural-extension-transition}
\end{equation}
On the quotient, every old boundary summand maps to the identically labelled
summand and every new-prime summand has zero preimage.  Thus ``old'' and
``new'' are labels in an explicit direct-sum map, not a claim of categorical
separation.

We next remove the last arbitrary generic stabilization.  For a generic
$x=((x_p),y)$ put
\begin{align}
 g_P(x)&=\operatorname{sgn}(y)\prod_{p\in P}p^{v_p(x_p)}\in\Gamma_P,
 \notag\\
 a_p(x)&=\operatorname{sgn}(y)
 \prod_{q\in P\setminus\{p\}}q^{-v_q(x_q)}
 p^{-v_p(x_p)}x_p,
 \notag\\
 t_P(x)&=\log|y|-\sum_{p\in P}v_p(x_p)\log p.
 \label{eq:coherence-generic-complete-coordinates}
\end{align}
This is the same retained coordinate as
\eqref{eq:multiplace-generic-coordinates}, with the unit
$u_p=p^{-v_p(x_p)}x_p$ written inside the formula.  The map
\begin{equation}
 \Omega_P:W_{\varnothing,P}\xrightarrow{\ \cong\ }
 \Gamma_P\times K_P\times\mathbb R,
 \qquad x\longmapsto(g_P(x),(a_p(x))_p,t_P(x))
 \label{eq:coherence-generic-homeomorphism}
\end{equation}
has inverse
\begin{equation}
 (g,a,t)\longmapsto g\cdot((a_p)_{p\in P},e^t),
 \label{eq:coherence-generic-homeomorphism-inverse}
\end{equation}
and $\delta\in\Gamma_P$ acts by
$(g,a,t)\mapsto(\delta g,a,t)$.  The same valuation-sector argument as in
Theorem~\ref{thm:coherence-sector-homeomorphism} proves the homeomorphism.
Consequently
\begin{equation}
 \mathfrak I_P\cong
 \mathcal K(\ell^2\Gamma_P)\otimes C(K_P)\otimes C_0(\mathbb R)
 \label{eq:coherence-generic-crossed-product-coordinate}
\end{equation}
by the exact matrix formula
\[
 (\mathbf1_{\{g\}}f)u_{gh^{-1}}\longmapsto e_{g,h}\otimes f.
\]

Let
\begin{equation}
 \mathcal H_{\mathbb Q}=\ell^2(\mathbb Q^\times)\otimes H_0,
 \qquad
 J_{\mathbb Q}=C_0(\mathbb R)\otimes\mathcal K(\mathcal H_{\mathbb Q}).
 \label{eq:coherence-universal-ideal}
\end{equation}
The basis unitary
\begin{equation}
 \ell^2(\Gamma_P)\otimes\ell^2(\mathcal L_P)
 \xrightarrow{\ \cong\ }\ell^2(\mathbb Q^\times),
 \qquad \delta_\gamma\otimes\delta_\ell\longmapsto\delta_{\gamma\ell}
 \label{eq:coherence-rational-basis-unitary}
\end{equation}
turns \eqref{eq:coherence-generic-crossed-product-coordinate} into the
simultaneous isomorphism
\begin{equation}
 \widetilde\chi_P:\widetilde{\mathfrak I}_P
 \xrightarrow{\ \cong\ } C(K_P)\otimes J_{\mathbb Q}.
 \label{eq:coherence-simultaneous-generic-coordinate}
\end{equation}
No unspecified bijection between separable Hilbert spaces occurs here: every
basis vector is labelled by the displayed rational and the retained $H_0$
index.

Let
\begin{equation}
 r_{P,P'}:K_{P'}=K_P\times K_R\longrightarrow K_P
 \label{eq:coherence-unit-projection}
\end{equation}
be projection.  It is onto, its fibre over every point is $K_R$, and its
kernel as a compact-group homomorphism is $\{1\}_{K_P}\times K_R$.  Direct
substitution of
\eqref{eq:coherence-old-coordinate}--
\eqref{eq:coherence-new-unit-coordinate} into
\eqref{eq:coherence-generic-complete-coordinates} proves the exact ideal
transition formula
\begin{equation}
 \widetilde\chi_{P'}\Psi^I_{P,P'}\widetilde\chi_P^{-1}
 =r_{P,P'}^*\otimes\operatorname{id}_{J_{\mathbb Q}}.
 \label{eq:coherence-exact-ideal-transition}
\end{equation}
Indeed, the $P$ unit coordinates become the $P$ coordinates of $K_{P'}$;
after the residual $\Gamma_P$ representative is fixed, the $b_r$ become the
$R$ coordinates.  On Hilbert-space matrix units, both sides use the same
identity
\[
 \Gamma_P\times\Gamma_R^+\times\mathcal L_{P'}
 =\Gamma_{P'}\times\mathcal L_{P'}=\mathbb Q^\times.
\]

Equation \eqref{eq:coherence-exact-ideal-transition} is nondegenerate and has
the unique strict multiplier extension.  In the standard multiplier-section
coordinate, for a bounded strictly continuous
$M:K_P\to M(J_{\mathbb Q})$, it is
\begin{align}
 (\overline\Psi^I_{P,P'}M)(c_P,c_R)&=M(c_P),
 \label{eq:coherence-ideal-multiplier-map}\\
 (\dot\Psi^I_{P,P'}[M])(c_P,c_R)&=[M(c_P)]
 \quad\text{in }Q(C(K_{P'})\otimes J_{\mathbb Q}).
 \label{eq:coherence-ideal-corona-map}
\end{align}
The second equation is a formula on multiplier representatives; it does not
assert that every corona class has a globally chosen continuous multiplier
representative.

Define the constant-unit embedding in these simultaneous coordinates by
\begin{equation}
 \widetilde\jmath_P:J_{\mathbb Q}\longrightarrow
 \widetilde{\mathfrak I}_P,
 \qquad
 \widetilde\chi_P(\widetilde\jmath_P(b))=1_{K_P}\otimes b.
 \label{eq:coherence-constant-ideal-embedding}
\end{equation}
Then, before and after passing to multipliers or coronas,
\begin{equation}
 \Psi^I_{P,P'}\widetilde\jmath_P=\widetilde\jmath_{P'},
 \qquad
 \overline\Psi^I_{P,P'}\overline{\widetilde\jmath}_P
 =\overline{\widetilde\jmath}_{P'},
 \qquad
 \dot\Psi^I_{P,P'}\dot{\widetilde\jmath}_P
 =\dot{\widetilde\jmath}_{P'}.
 \label{eq:coherence-constant-multiplier-corona-squares}
\end{equation}
The image of the first map in
\eqref{eq:coherence-exact-ideal-transition} consists exactly of the sections
constant on each $K_R$ fibre.  All nonconstant new-unit sections remain in
the target; none has been declared equal to an old section.

\begin{proposition}[Triple inclusions compose on the nose]
\label{prop:coherence-triple-inclusions}
If $P\subset P'\subset P''$, then, with the tensor factors ordered by the
increasing numerical order of their prime labels,
\begin{equation}
 \Psi_{P',P''}\Psi_{P,P'}=\Psi_{P,P''}
 \label{eq:coherence-triple-middle-composition}
\end{equation}
and the same identities hold for the ideal and quotient maps and their ideal
multiplier and corona extensions.
\end{proposition}

\begin{proof}
Put $R=P'\setminus P$ and $T=P''\setminus P'$.  For
$x\in U_{P,P''}$, direct extraction gives
\[
 h_{R\sqcup T}(x)=h_R(x)h_T(x).
\]
First extracting $T$ divides every remaining coordinate, including every
$R$ coordinate, by $h_T(x)$.  Extracting $R$ next therefore gives
\[
 (h_Rh_T)^{-1}x_q\quad(q\in P),\qquad
 (h_Rh_T)^{-1}x_s\quad(s\in R\sqcup T),
\]
which are exactly the direct old and unit coordinates.  Thus the two
homeomorphisms agree after the displayed reassociation of the labelled
factors.

On tail bases, unique factorization gives the same equality
\[
 \mathcal L_P=\Gamma_R^+\times\Gamma_T^+\times\mathcal L_{P''}.
\]
Hence both compositions take
$e_{\rho\tau\ell,\,\sigma\upsilon m}$ to
$e_{\rho,\sigma}\otimes e_{\tau,\upsilon}\otimes e_{\ell,m}$.
Coefficient pullbacks and extension by zero compose literally.  Formula
\eqref{eq:coherence-sector-matrix-units} therefore proves
\eqref{eq:coherence-triple-middle-composition} on a dense algebraic span,
and continuity proves it everywhere.  Restriction, quotient, and the unique
strict extension of the nondegenerate ideal map preserve the equality.

This is the concrete instance of the standard induction-in-stages principle:
composition of groupoid correspondences becomes the internal tensor product,
and staged induction equals direct induction
\cite[Section~8.4.3, p.~35]{li2024groupoidnotes}.  Li's notes place that
principle in the groupoid framework developed from Williams's text
\cite{williams2019toolkit}; the equality used here has nevertheless been
proved above on the actual coordinates and matrix units.
\end{proof}

\subsubsection{Boundary carriers and the connecting-map square}

For $p\in P$, recall
$Y_{p,P}=\prod_{q\in P\setminus\{p\}}\mathbb Z_q^\times$.  There are literal
factorizations and projections
\begin{equation}
 Y_{p,P'}=Y_{p,P}\times K_R,
 \qquad
 \pi_{p;P,P'}:Y_{p,P'}\longrightarrow Y_{p,P}.
 \label{eq:coherence-finite-boundary-unit-projection}
\end{equation}
They intertwine translation by
$g_{p,P'}=(g_{p,P},(p)_{r\in R})$ and by $g_{p,P}$.  Therefore
\begin{equation}
 \pi_{p;P,P'}(H_{p,P'})=H_{p,P}
 \label{eq:coherence-procyclic-closure-projection}
\end{equation}
and there is a surjection
\begin{equation}
 \overline\pi_{p;P,P'}:
 Y_{p,P'}/H_{p,P'}\longrightarrow Y_{p,P}/H_{p,P}.
 \label{eq:coherence-finite-orbit-quotient-projection}
\end{equation}
At the mapping-torus level the map is
\begin{equation}
 B_{p,P'}\longrightarrow B_{p,P},\qquad
 [(y,k_R),s]\longmapsto[y,s].
 \label{eq:coherence-finite-mapping-torus-projection}
\end{equation}
It is well defined because both mapping-torus actions multiply every retained
unit by $p$ and translate $s$ by $\log p$.

Under the complete K-groups of
Lemma~\ref{lem:unit-fibre-mapping-torus-k-groups}, the old finite summand of
$(\Psi^D_{P,P'})_*$ is
\begin{align}
 h&\longmapsto h\circ\overline\pi_{p;P,P'}
 &&\text{on }K_0,
 \label{eq:coherence-finite-K0-transition}\\
 [f]&\longmapsto[f\circ\pi_{p;P,P'}]
 &&\text{on }K_1.
 \label{eq:coherence-finite-K1-transition}
\end{align}
The second formula is well defined because pullback commutes with
$1-\sigma_{p,P}$.  No injectivity assertion about the coinvariant map is
being inserted without proof.  The $K_0$ map is injective because
\eqref{eq:coherence-finite-orbit-quotient-projection} is onto.

For the archimedean summand use the choice-free orbit carrier
\begin{equation}
 \overline K_P=K_P/\{c\sim-c\},\qquad
 q_P^\pm:K_P\longrightarrow\overline K_P.
 \label{eq:coherence-arch-orbit-carrier}
\end{equation}
The projection $r_{P,P'}$ is sign-equivariant and induces
\begin{equation}
 \overline r_{P,P'}:\overline K_{P'}\longrightarrow\overline K_P,
 \qquad[(c_P,c_R)]\longmapsto[c_P].
 \label{eq:coherence-arch-orbit-projection}
\end{equation}
Under the canonical Morita coordinate
$K_0(\mathfrak D_{\infty,P})=C(\overline K_P,\mathbb Z)$, its transition is
$\overline r_{P,P'}^*$.  Choosing a clopen half $H_P^+$ recovers the earlier
$C(H_P^+,\mathbb Z)$ display, but no compatibility of independently chosen
halves is required here.

Define the complete old-boundary $K_0$ transition by
\begin{equation}
 T^0_{P,P'}((h_p)_{p\in P},h_\infty)
 =\bigl((h_p\circ\overline\pi_{p;P,P'})_{p\in P},
        (0)_{r\in R},
        h_\infty\circ\overline r_{P,P'}\bigr).
 \label{eq:coherence-global-boundary-K0-transition}
\end{equation}
The complete natural connecting map, in the same choice-free archimedean
coordinate, is
\begin{equation}
 \partial_P^{(1)}((h_p)_p,h_\infty)
 =-\sum_{p\in P}\iota_{p,P}(h_p)-h_\infty\circ q_P^\pm.
 \label{eq:coherence-choice-free-global-boundary}
\end{equation}
This is equation
\eqref{eq:global-full-unit-fibre-boundary-map} transported through the
homeomorphism $H_P^+\cong\overline K_P$.

\begin{theorem}[Exact finite-set connecting square]
\label{thm:coherence-connecting-square}
For every $P\subset P'$,
\begin{equation}
 r_{P,P'}^*\partial_P^{(1)}
 =\partial_{P'}^{(1)}T^0_{P,P'}.
 \label{eq:coherence-connecting-square}
\end{equation}
Equivalently, this is the $K$-theory naturality square induced by
\eqref{eq:coherence-natural-extension-transition}.  It retains all finite
invariant functions, the complete archimedean sign orbit, and zero-labelled
new boundary summands.
\end{theorem}

\begin{proof}
The definitions give the two literal identities
\begin{align}
 r_{P,P'}^*\iota_{p,P}(h)
 &=\iota_{p,P'}(h\circ\overline\pi_{p;P,P'}),
 \label{eq:coherence-finite-boundary-square}\\
 r_{P,P'}^*(h_\infty\circ q_P^\pm)
 &=(h_\infty\circ\overline r_{P,P'})\circ q_{P'}^\pm.
 \label{eq:coherence-arch-boundary-square}
\end{align}
For the first, both sides evaluate $h$ after dropping the $p$ coordinate and
then forgetting every new-prime unit coordinate.  For the second, both sides
send $(c_P,c_R)$ to the sign orbit $[c_P]$.  Substitute
\eqref{eq:coherence-finite-boundary-square}--
\eqref{eq:coherence-arch-boundary-square} into
\eqref{eq:coherence-choice-free-global-boundary}; every new-prime boundary
coordinate in \eqref{eq:coherence-global-boundary-K0-transition} is zero.
The displayed equality follows with its minus sign unchanged.  It also follows
from the published naturality theorem for the already constructed extension
morphism \cite[Definition~21.1.1 and Example~21.1.2(a)]{blackadar1986}, but
the coordinate calculation proves exactly which functions that abstract
square contains.
\end{proof}

On $K_1(\mathfrak D_P)$, the transition is the direct sum of
\eqref{eq:coherence-finite-K1-transition} in the old prime coordinates and
zero in the new coordinates.  The six-term sequence then supplies the induced
middle-algebra maps.  No splitting of the middle $K_1$ extension and no
freeness claim for the profinite coinvariant groups is selected.

\subsubsection{A coherent stable-to-natural family}

We now apply the natural transition maps to the stable balancing source.
Tensor the stable source extension
\eqref{eq:global-place-stable-extension} by
$\mathcal K(\ell^2\mathbb Q^\times)$, reorder the two explicitly labelled
Hilbert factors, and denote it by
\begin{equation}
 0\longrightarrow J_{\mathbb Q}\longrightarrow
 \widetilde{\mathcal E}_P\longrightarrow
 \widetilde{\mathcal A}_P\longrightarrow0.
 \label{eq:coherence-tail-stabilized-source-extension}
\end{equation}
The source inclusion for $P\subset P'$ is
\begin{equation}
 \widetilde\iota^{\mathcal E}_{P,P'}(a,m)
 =\bigl(\iota_{P,P'}(a),m\bigr),
 \label{eq:coherence-source-extension-transition}
\end{equation}
where $\iota_{P,P'}$ inserts zero in every new-prime quotient coordinate.
It is the identity on $J_{\mathbb Q}$ and uses the same fixed place-indexed
corner isometries at every finite stage.

There is a particularly economical way to make the target edge maps
coherent.  For a finite place $p$, fix once the base edge morphism at
$P_p=\{p\}$ supplied by
Lemma~\ref{lem:compatible-stabilized-edge-morphisms}, after the displayed tail
stabilization.  For the archimedean place use $P_\infty=\varnothing$.  Only
for this base object, extend the preceding definitions by
\[
 K_\varnothing=\{1\},\qquad
 \Gamma_\varnothing=\{\pm1\},\qquad
 \mathbb A_{S_\varnothing}=\mathbb R,
 \]
and take the literal reflection extension
$C_0(\mathbb R)\rtimes\{\pm1\}$ with generic open set
$\mathbb R^\times$ and boundary $\{0\}$.  Equations
\eqref{eq:coherence-complementary-prime-group}--
\eqref{eq:coherence-natural-middle-transition-expanded} then apply with
$R=P$: the old-boundary sector is precisely the archimedean edge.  Its
generic, middle, and boundary base maps are the identity before the same tail
stabilization.  For every
$P\supset P_v$ define
\begin{equation}
 (\widetilde\jmath_P,
  \widetilde\lambda_{v,P},
  \widetilde\alpha_{v,P})
 =
 (\Psi^I_{P_v,P},\Psi_{P_v,P},\Psi^D_{P_v,P})
 \circ
 (\widetilde\jmath_{P_v},
  \widetilde\lambda_{v,P_v},
  \widetilde\alpha_{v,P_v}).
 \label{eq:coherence-placewise-edge-definition}
\end{equation}
The ideal component agrees with
\eqref{eq:coherence-constant-ideal-embedding} by
\eqref{eq:coherence-constant-multiplier-corona-squares}.  Proposition
\ref{prop:coherence-triple-inclusions} makes this definition independent of
an intermediate finite set and gives
\begin{align}
 \Psi_{P,P'}\widetilde\lambda_{v,P}
 &=\widetilde\lambda_{v,P'},
 \label{eq:coherence-placewise-middle-square}\\
 \Psi^D_{P,P'}\widetilde\alpha_{v,P}
 &=\widetilde\alpha_{v,P'}
 \qquad(v\in P\sqcup\{\infty\}).
 \label{eq:coherence-placewise-quotient-square}
\end{align}

Let $S_{\nu(v)}$ be the fixed isometries of
\eqref{eq:global-place-cuntz-family}, now acting on the final $H_0$ factor.
Because every $\Psi_{P,P'}$ is the identity on that factor, its corner
operation commutes exactly with transition.  Form the quotient sum
\begin{equation}
 \widetilde\alpha_P((a_v)_v)
 =\sum_{v\in P\sqcup\{\infty\}}
 (1\otimes S_{\nu(v)})
 \widetilde\alpha_{v,P}(a_v)
 (1\otimes S_{\nu(v)})^*.
 \label{eq:coherence-global-quotient-embedding}
\end{equation}
Its Busby equation is
\begin{equation}
 \widetilde\tau_P^{\mathrm{nat}}\widetilde\alpha_P
 =\dot{\widetilde\jmath}_P\widetilde\tau_P.
 \label{eq:coherence-global-busby-equation}
\end{equation}
It follows place by place from the base Busby equations, then from
\eqref{eq:coherence-constant-multiplier-corona-squares}, and finally from the
orthogonal corner sum.  Define in the two Busby pullbacks
\begin{equation}
 \widetilde\Phi_P(a,m)
 =\bigl(\widetilde\alpha_P(a),
        \overline{\widetilde\jmath}_P(m)\bigr).
 \label{eq:coherence-global-middle-embedding}
\end{equation}

\begin{theorem}[Simultaneous finite-set coherence of the stable embedding]
\label{thm:coherence-stable-to-natural-family}
The maps \eqref{eq:coherence-global-middle-embedding} are injective morphisms
of the tail-stabilized source extensions into the tail-stabilized natural
extensions.  For every inclusion $P\subset P'$ of nonempty finite prime sets,
\begin{equation}
 \Psi_{P,P'}\widetilde\Phi_P
 =\widetilde\Phi_{P'}\widetilde\iota^{\mathcal E}_{P,P'}.
 \label{eq:coherence-stable-to-natural-family}
\end{equation}
The same square commutes on ideals, quotients, ideal multipliers, ideal
coronas, and connecting maps.  For three nested finite sets, the square is
strictly associative by
\eqref{eq:coherence-triple-middle-composition}.
\end{theorem}

\begin{proof}
Equation \eqref{eq:coherence-global-busby-equation} puts
\eqref{eq:coherence-global-middle-embedding} in the natural Busby pullback.
The constant ideal map and every orthogonal quotient corner are injective.
The same two-outside-arrows argument as in
Theorem~\ref{thm:stable-to-natural-extension-embedding} proves injectivity of
the middle map.

For $a=(a_v)_{v\in P\sqcup\{\infty\}}$, equations
\eqref{eq:coherence-placewise-quotient-square} and
\eqref{eq:coherence-global-quotient-embedding}, together with identity on the
$H_0$ corner factor, give
\[
 \Psi^D_{P,P'}\widetilde\alpha_P(a)
 =\widetilde\alpha_{P'}(\iota_{P,P'}a).
\]
The new quotient coordinates on the right are zero.  On multiplier
coordinates, equation
\eqref{eq:coherence-constant-multiplier-corona-squares} gives
\[
 \overline\Psi^I_{P,P'}
 \overline{\widetilde\jmath}_P(m)
 =\overline{\widetilde\jmath}_{P'}(m).
\]
The two coordinates in the natural Busby pullback are therefore equal, which
proves \eqref{eq:coherence-stable-to-natural-family}.  Restriction and quotient
give the ideal and boundary squares; strict extension gives the ideal
multiplier and corona squares; and
Theorem~\ref{thm:coherence-connecting-square} gives the complete connecting
map.  Triple associativity is already proved on the middle algebras and hence
persists on every induced component.
\end{proof}

This theorem constructs a coherent replacement for the independently chosen
coordinates in Subsection~\ref{subsec:stable-to-natural-extension}.  It does
not assert that an arbitrary prior family $(\chi_P)_P$ was secretly coherent.
For any such family, the exact unresolved datum for an inclusion is the pair
of homomorphisms
\begin{equation}
 \chi_{P'}\Psi^I_{P,P'}
 \quad\text{and}\quad
 (r_{P,P'}^*\otimes\operatorname{id})\chi_P.
 \label{eq:coherence-arbitrary-choice-defect}
\end{equation}
Coherence means equality of this pair after placing both maps in the same
declared stabilization.  Formula
\eqref{eq:coherence-simultaneous-generic-coordinate} makes the pair equal by
construction.  No quotient by a vague ``unitary equivalence'' is taken.

\subsubsection{The exact middle-multiplier obstruction}

For $P\subsetneq P'$, let
\begin{equation}
 \widetilde{\mathfrak O}_{P,P'}=
 (C_0(U_{P,P'})\rtimes\Gamma_{P'})
 \otimes\mathcal K(\ell^2\mathcal L_{P'})\otimes\mathcal K(H_0).
 \label{eq:coherence-old-sector-ideal}
\end{equation}
The map $\Psi_{P,P'}$ is nondegenerate as a map into
$\widetilde{\mathfrak O}_{P,P'}$ and therefore has a unique strict extension
\begin{equation}
 M(\widetilde{\mathfrak N}_P)longrightarrow
 M(\widetilde{\mathfrak O}_{P,P'}).
 \label{eq:coherence-sector-multiplier-map}
\end{equation}
The sector is an essential ideal in
$\widetilde{\mathfrak N}_{P'}$: it contains the dense generic coefficient
ideal, and the Fourier-coefficient argument of
Lemma~\ref{lem:natural-one-zero-essential-pullback} applies verbatim.
It is nevertheless proper.  For every $r\in R$, the nonzero new boundary
summand $\widetilde{\mathfrak D}_{r,P'}$ is annihilated by the sector quotient.

Consequently $\Psi_{P,P'}$ is degenerate as a map into the full
$\widetilde{\mathfrak N}_{P'}$.  If it had a unital strict extension
$M(\widetilde{\mathfrak N}_P)\to M(\widetilde{\mathfrak N}_{P'})$, the image
of an approximate unit would converge strictly to one.  Applying the quotient
onto $\widetilde{\mathfrak D}_{r,P'}$ would make the identically zero image
converge strictly to its nonzero identity multiplier, a contradiction.
Thus the canonical middle multiplier target is
$M(\widetilde{\mathfrak O}_{P,P'})$, not
$M(\widetilde{\mathfrak N}_{P'})$.  Essentiality embeds
$M(\widetilde{\mathfrak N}_{P'})$ into
$M(\widetilde{\mathfrak O}_{P,P'})$ in the opposite direction; reversing
that arrow would be an unsupported categorical assertion.

This obstruction does not damage the Busby construction.  The generic ideal
map is nondegenerate into the complete generic ideal and has exactly the
multiplier and corona maps
\eqref{eq:coherence-ideal-multiplier-map}--
\eqref{eq:coherence-ideal-corona-map}.  Those are the maps used in
\eqref{eq:coherence-global-busby-equation} and in the pullback formula
\eqref{eq:coherence-global-middle-embedding}.

\paragraph{Exact boundary of the result.}
The finite-set system now has explicit maps on compact unit carriers, generic
ideals, finite and archimedean boundary quotients, tail Hilbert spaces,
orthogonal corners, Busby pullbacks, ideal multipliers, ideal coronas, and
connecting maps.  All triple-inclusion squares commute in the declared basis.
The proper old-sector ideal proves why there is no unital full-middle
multiplier map when a prime is genuinely added.  The theorem is an existence
and construction result for a coherent coordinate family; it is not a
choice-independence theorem for every possible stabilization.  It does not
select a splitting of any coinvariant or middle $K_1$ sequence, identify a
new-prime unit fibre with a point, or imply an endpoint, individual-zero,
critical-line, or global zeta statement.


\subsection{Candidate interfaces retained without overclaiming}

The following broader source interfaces remain useful research leads.  The
table separates the bounded theorems now proved from the additional evidence
that would be required for any stronger interface claim.

\begin{center}
\begin{tabular}{p{0.25\textwidth}p{0.31\textwidth}p{0.34\textwidth}}
\toprule
Source layer & Precisely retained datum & Missing typed evidence before admission\\
\midrule
 Global finite-prime assembly & The local boundary is $-1$; the natural
 unit-retaining multi-place algebra is Definition
 \ref{def:unstabilized-one-zero-algebra}; Theorem
 \ref{thm:stable-to-natural-extension-embedding} lifts the stable codiagonal
 square to an injective extension morphism; and Theorem
 \ref{thm:full-natural-one-zero-k-theory} computes the full nonconstant
 finite invariants, finite coinvariants, archimedean sign-even image, natural
 connecting map, $K_0$ kernel, and $K_1$ extension & Any stronger
 surjectivity, isomorphism, or Morita-equivalence claim must account for the
 explicit odd-difference quotient $\mathcal Q_P$, the finite-level
 coinvariant multiplicities, and coherence under enlarging $P$.
 Proposition \ref{prop:minimal-symmetric-sign-enrichment} retains the sign
 difference in a separate split carrier; it does not identify that carrier
 with the natural boundary algebra\\
Archimedean APS/global splice & The reflection crossed product, its matrix
 model, quotient, two boundary generators, map $\partial_\infty(a,b)=-(a+b)$,
 the whole-line and two compact-interval index realizations, and the stable
 multi-place codiagonal and the natural arithmetic representative are proved
 separately & A specified APS domain, endpoint projectors, adjoint domain,
 and an exact morphism from that operator realization to the extension data\\
Orbifold labels & The conditional $(C,\rho)$ label set and the $S_3$ count $(1,1,2,3,3,2,2,2)$ & The exact VOA/orbifold hypotheses and a content-read primary source for the claimed fusion setting\\
Leech/moonshine & The formal $A_2$ direct-sum rank count and the displayed theta identity as candidate data & A human-source chain for the $12A_2$ neighbour construction and the stated simple-current identifications\\
Spectral/F-theory layers & Rank-one, ADE, Mordell--Weil, and BSD expressions as separately typed analogies or sourced dictionaries & Exact operators, bases, arithmetic fibration, and maps; no common object follows from shared notation\\
CUE/Epstein bridge & The four listed structural coincidences & An explicit transform-level map between compact pushforward measures and discrete lattice-shell measures\\
\bottomrule
\end{tabular}
\end{center}

Consequently, the local source's broad ``primitive defect realization theorem''
is replaced by a collection of theorems above plus an open-interface ledger.
No endpoint statement, individual-zero map, critical-line conclusion, or
Riemann-zeta zero result is inferred from the finite algebraic coincidences.
